% !TEX program = xelatex
\documentclass[11pt,openany,fontset=fandol]{ctexbook}
\usepackage[paperwidth=176mm,paperheight=250mm,inner=21mm,outer=18mm,top=20mm,bottom=20mm,headheight=14pt,headsep=6mm,footskip=10mm]{geometry}
\usepackage{amsmath,amssymb,amsthm,bm,mathtools,booktabs,array,enumitem,float}
\usepackage{tikz,fancyhdr,hyperref}
\usetikzlibrary{arrows.meta,calc,patterns,angles,quotes,decorations.markings,decorations.pathmorphing}
\theoremstyle{definition}\newtheorem{example}{例}[chapter]\newtheorem{drill}{练习}[chapter]
\newenvironment{solution}{\par\smallskip\noindent\textbf{解}\quad\ignorespaces}{\par\medskip}
\newenvironment{drillsolution}[1]{\subsubsection*{练习\ref{#1}详解}\phantomsection\label{sol:#1}\ignorespaces}{\par\medskip}
\newenvironment{question}{\par\medskip\noindent\begin{minipage}{\linewidth}}{\end{minipage}\par\medskip}
\newcommand{\dd}{\mathop{}\!\mathrm d}
\newcommand{\e}{\mathrm e}
\newcommand{\fig}[3]{\begin{figure}[H]\centering #3\caption{#2}\label{#1}\end{figure}}
\newcommand{\answerlink}[1]{\par\smallskip\noindent\hyperref[sol:#1]{详解}}
\tikzset{axis/.style={->,>=Stealth,line width=.45pt},world/.style={line width=.8pt},flow/.style={->,>=Stealth,line width=.7pt},ray/.style={->,>=Stealth,line width=.7pt},guide/.style={dashed,line width=.35pt},aux/.style={dashed,line width=.35pt},dim/.style={<->,>=Stealth,line width=.35pt},every node/.style={font=\small},force/.style={->,>=Stealth,very thick},gas/.style={fill=blue!7}}
% Common visual style for all first-round physics lecture projects.
% Each project includes an identical local copy of this file.
\geometry{paperwidth=176mm,paperheight=250mm,inner=21mm,outer=18mm,top=20mm,bottom=20mm,headheight=14pt,headsep=6mm,footskip=10mm}
\hypersetup{hidelinks}
\setlength{\parindent}{2em}\setlength{\parskip}{3pt}\setlength{\emergencystretch}{2em}
\linespread{1.18}\raggedbottom\setcounter{tocdepth}{1}
\widowpenalty=10000\clubpenalty=10000\displaywidowpenalty=10000
\setlength{\intextsep}{8pt plus1pt minus1pt}
\makeatletter
\g@addto@macro\normalsize{\setlength{\abovedisplayskip}{8pt plus2pt minus2pt}\setlength{\belowdisplayskip}{8pt plus2pt minus2pt}\setlength{\abovedisplayshortskip}{0pt plus2pt}\setlength{\belowdisplayshortskip}{6pt plus2pt minus2pt}}
\makeatother
\providecommand{\Needspace}[1]{}
\newdimen\LectureAvailableSpace
\renewcommand{\Needspace}[1]{\par\penalty-100\begingroup\LectureAvailableSpace=\pagegoal\advance\LectureAvailableSpace by-\pagetotal\ifdim\LectureAvailableSpace<#1\relax\ifdim\pagetotal>0pt\newpage\fi\fi\endgroup}
\renewenvironment{solution}{\par\smallskip\Needspace{3\baselineskip}\noindent\textbf{解}\quad\ignorespaces}{\par\medskip}
\ctexset{chapter={name={第,讲},number=\chinese{chapter},format=\LARGE\bfseries\centering,beforeskip=16pt,afterskip=22pt},section={format=\Large\bfseries},subsection={format=\large\bfseries},subsubsection={format=\normalsize\bfseries,beforeskip=10pt,afterskip=4pt}}
\setlist[enumerate]{leftmargin=2em,itemsep=2pt,topsep=3pt}
\pagestyle{fancy}\fancyhf{}\fancyhead[LE,RO]{\small\thepage}\fancyhead[LO]{\small 物理竞赛一轮讲义}
\renewcommand{\headrulewidth}{0pt}
\fancypagestyle{plain}{\fancyhf{}\fancyfoot[C]{\small\thepage}\renewcommand{\headrulewidth}{0pt}}
\renewenvironment{drillsolution}[1]{\par\Needspace{4\baselineskip}\subsubsection*{练习\ref{#1}详解}\phantomsection\label{sol:#1}\nobreak\noindent\textbf{解}\quad\ignorespaces}{\par\medskip}
\ifdefined\exercisesolution
\renewenvironment{exercisesolution}[1]{\par\Needspace{4\baselineskip}\subsubsection*{习题\ref{#1}详解}\phantomsection\label{sol:#1}\nobreak\noindent\textbf{解}\quad\ignorespaces}{\par\medskip}
\fi
\newcommand{\LectureAnswerSection}{\par\section*{练习解答}\addcontentsline{toc}{section}{练习解答}}
\newcommand{\LectureHeader}[2]{\fancyhead[RE]{\small #1\quad #2}}
\newcommand{\LectureCover}[3]{\begin{titlepage}\centering\vspace*{32mm}{\Huge\bfseries #1\par}\vspace{12mm}{\Large 物理竞赛一轮讲义\par}\vspace{18mm}{\LARGE #2\\[5mm]#3\par}\vfill{\large 2026\par}\end{titlepage}}

\LectureHeader{第三讲}{极值、近似与曲线几何}
\begin{document}\frontmatter\LectureCover{微积分}{第三讲}{极值、近似与曲线几何}\tableofcontents\mainmatter
\setcounter{chapter}{2}\chapter{极值、近似与曲线几何}
\section{极值、最值与平衡}
\subsection{导数符号比零点方程包含更多信息}
在区间内部，若 $f'>0$，函数随自变量增大而增大；若 $f'<0$，则减小。可导函数在内点取得局部极值时，必有 $f'=0$，但这个条件只是候选条件。例如 $x^3$ 在零点导数为零，却始终递增，因此零点不是极值点。

判断候选点应看两侧导数：由正变负是极大，由负变正是极小；若导数不变号，通常不是极值。二阶导数判据是这一规律的局部形式：若 $f'(x_0)=0$，且 $f''(x_0)>0$，则附近向上弯曲，$x_0$ 为极小；若 $f''(x_0)<0$ 则为极大。二阶导数为零时必须继续分析，不能直接下结论。

闭区间上的连续函数必有最大值、最小值。求全局最值时，应把内部驻点、内部不可导点与端点一起比较。端点未必满足 $f'=0$，例如 $f(x)=x$ 在 $[0,1]$ 上的最大值位于 $1$。开区间则可能只有上确界或下确界，而没有实际取得该值的位置；定义域是解题步骤的一部分。

导数判单调的依据是有限变化也能由某个中间位置的变化率控制。若 $f$ 在 $[a,b]$ 上连续、内部可导，中值定理给出 $f(b)-f(a)=f'(\xi)(b-a)$。所以整个区间内导数为正时，任意从左到右的有限变化都为正。仅知道某一点导数为零，没有提供两侧有限变化的符号，自然不足以判断极值。

实际建模时先确定允许区间，再寻找候选位置。例如绳长、折射角、液面高度有各自边界，数学驻点可能落在不可实现区域。对含绝对值或接触切换的关系，还要把表达式发生变化的位置列入比较。得到候选点后代回函数值，而不是比较导数值，才能回答题目所问的最大量或最小量。

\fig{fv:c3:stationary}{同为水平切线，导数是否变号决定极值性质}{\begin{tikzpicture}[x=.9cm,y=.75cm]
\foreach \x/\lab in {0/{$x^2$},4.5/{$x^3$}}{\begin{scope}[xshift=\x cm]
\draw[->](-1.7,0)--(1.9,0)node[right]{$x$};\draw[->](0,-1.5)--(0,2.2)node[above]{$f$};\node at(1.3,1.9){\lab};\end{scope}}
\draw[blue,thick,domain=-1.35:1.35,samples=60]plot(\x,{\x*\x});
\draw[red,thick,domain=-1.1:1.1,samples=60]plot({4.5+\x},{\x*\x*\x});
\node at(-.9,-.6){$f'<0$};\node at(.9,-.6){$f'>0$};
\node at(3.5,-1.7){$f'>0$};\node at(5.4,-.6){$f'>0$};
\end{tikzpicture}}
图\ref{fv:c3:stationary}左图从原点左侧减小、右侧增大，原点确为极小点；右图经过原点仍递增，水平切线只表示瞬时斜率为零。因而解 $f'=0$是找候选点，读取两侧变化才完成判别。若自变量还限于闭区间，例如 $x^3$在 $[-1,1]$内求最值，最大、最小分别在端点，原点的驻点性质不能替代端点比较。


\subsection{势能导数与稳定性}
对一个广义坐标 $q$，保守力的虚功满足 $Q\,\dd q=-\dd U$，故 $Q=-U'(q)$。若 $q$ 是长度，$Q$ 为相应方向的力；若 $q$ 是角度，$Q$ 是力矩。不能不看坐标意义就把势能的一阶导数都称作同一种力。

平衡满足 $U'(q_0)=0$。作小位移 $\eta=q-q_0$，若
\[
 U(q)=U(q_0)+\tfrac12K\eta^2+\cdots,
\]
则恢复作用为 $Q=-K\eta+\cdots$。$K>0$ 时扰动受到恢复，称稳定平衡；$K<0$ 时扰动被放大，称不稳定平衡。若 $K=0$，应看更高阶：$+\eta^4$ 仍给出稳定平衡，$-\eta^4$ 不稳定，$\eta^3$ 也不能给出局部势阱。若一个区间内势能完全恒定，则在该坐标方向为随遇平衡。

稳定性还依赖允许的扰动。只准沿杆轴改变两球距离时的稳定，不等于对所有空间扰动都稳定。题设中的约束必须保留到判据中。

\fig{c3k-potential}{势能在平衡点两侧的变化决定扰动后的趋势}{\begin{tikzpicture}[x=1cm,y=.9cm]
\draw[axis](-2.2,0)--(2.5,0)node[right]{$q-q_0$};
\draw[axis](0,-1.9)--(0,2.1)node[above]{$U-U_0$};
\draw[world,blue,domain=-1.5:1.5,samples=40]plot(\x,{.65*\x*\x});
\draw[world,red,domain=-1.5:1.5,samples=40]plot(\x,{-.65*\x*\x});
\node[right,blue]at(1.55,1.46){$K>0$};
\node[right,red]at(1.55,-1.46){$K<0$};
\node[below left]at(0,0){$O$};
\end{tikzpicture}}
图 \ref{c3k-potential} 的蓝色曲线在平衡点两侧上升，向任一侧偏移都要增加势能，负梯度就把物体推回；红色曲线则相反。这里判定的是静态平衡位置对微小扰动的响应，不是说稳定系统一定最终停在平衡点。没有耗散时，稳定平衡附近通常发生往复运动；存在耗散时，才可能逐渐回到平衡点。

当二阶系数为零，继续求导只是找最低非零项的手段。若 $U-U_0=a\eta^{2n}$、$a>0$，两侧势能仍上升，因而仍稳定；若最低非零项为奇次幂，一侧势能下降，便不是局部势阱。对只允许单侧扰动的边界平衡，不能照搬双侧判据，应在真正允许的方向上检查势能变化。


图\ref{c3deep:quartic}把二阶判据失效时的三种情况画在同一尺度下。三条曲线在原点都满足一阶、二阶导数为零，但蓝线 $+\eta^4$ 在两侧都上升，红线 $-\eta^4$ 在两侧都下降，虚线 $\eta^3$ 则一侧升、一侧降。对蓝线，力 $Q=-4a\eta^3$ 仍然指回平衡点，只是不再正比于位移；因此稳定不能自动推出简谐振动。
\fig{c3deep:quartic}{同样具有零二阶导数，最低非零项仍可导致不同稳定性}{\begin{tikzpicture}[x=1.8cm,y=1.2cm]
\draw[axis](-1.5,0)--(1.5,0)node[right]{$\eta$};\draw[axis](0,-1.9)--(0,1.9)node[above]{$U-U_0$};
\draw[blue,thick,domain=-1.1:1.1,samples=70]plot(\x,{(\x)^4});\draw[red,thick,domain=-1.1:1.1,samples=70]plot(\x,{-(\x)^4});\draw[gray,dashed,domain=-1.1:1.1,samples=70]plot(\x,{(\x)^3});
\node[blue,right]at(1.1,1.46){$+a\eta^4$};\node[red,right]at(1.1,-1.46){$-a\eta^4$};\node[gray,left]at(-1.1,-1.33){$b\eta^3$};
\draw[flow,blue](.9,.15)--(.3,.15);\draw[flow,blue](-.9,.15)--(-.3,.15);\node[below left]at(0,0){$0$};
\end{tikzpicture}}


\subsection{平衡坐标改变后，惯量也要改变}
势能曲线的最低点确定稳定位置，却不能单独给出振动频率。若系统用坐标 $q$描述，动能为 $T=J(q)\dot q^2/2$，在平衡点附近令 $q=q_0+\eta$，最低阶动能为 $J(q_0)\dot\eta^2/2$，势能为 $U_0+U''(q_0)\eta^2/2$。对保守运动，先对精确能量逐项求导：
\[
 \frac{\dd(T+U)}{\dd t}
 =\frac12J'(q)\dot q^3+J(q)\dot q\ddot q+U'(q)\dot q=0.
\]
在 $\dot q\ne0$ 的运动段除以 $\dot q$，再以连续性接到转向点，得到 $J(q)\ddot q+J'(q)\dot q^2/2+U'(q)=0$。令位移和速度均为一阶小量，$J'(q)\dot q^2/2$ 至少二阶；$U'(q_0)=0$ 且 $U'(q_0+\eta)=U''(q_0)\eta+O(\eta^2)$。保留一阶，线性方程就是
\[J(q_0)\ddot\eta+U''(q_0)\eta=0,\qquad
\omega^2=\frac{U''(q_0)}{J(q_0)}.\]
坐标是长度时 $J$通常是有效质量；坐标是角度时则是相应转动惯量。若把角坐标换成弧长 $s=R\eta$，势能二阶系数与惯量都会除以 $R^2$，比值不变，实际频率自然不依赖坐标名称。

图\ref{c3k-coordinate}表示同一个约束位移的两种计量。只改势能导数、却仍用原来的质量或惯量，会多出或少掉 $R^2$。对两球相对距离的扰动，也必须使用相对运动的有效质量，而非任取一个球的质量。
\fig{c3k-coordinate}{角位移与弧长位移描述同一局部运动}{\begin{tikzpicture}
\draw[world](0,0)--(2.2,0);\draw[world](0,0)--(30:2.2);
\draw[blue,very thick](2.2,0)arc(0:30:2.2);
\draw(.7,0)arc(0:30:.7)node[midway,right]{$\eta$};
\node[below]at(1.1,0){$R$};\node[right]at(2.2,.7){$s=R\eta$};
\node[below left]at(0,0){$O$};
\end{tikzpicture}}



\section{近似展开与误差阶数}
\subsection{泰勒公式的结构}
当函数在 $x_0$ 附近足够光滑时，
\begin{equation}\label{eq:c3taylor}
 f(x_0+h)=f(x_0)+f'(x_0)h+
 \frac{f''(x_0)}{2!}h^2+\cdots+
 \frac{f^{(n)}(x_0)}{n!}h^n+R_n.
\end{equation}
式\eqref{eq:c3taylor}的有限阶展开只需要局部光滑性。若可用拉格朗日余项，
\[
 R_n=\frac{f^{(n+1)}(\xi)}{(n+1)!}h^{n+1},
\]
便能从导数界估计误差。把有限阶展开进一步写成无穷级数，还需证明余项趋零；并非每个无限次可导函数都等于其泰勒级数。

记 $R=O(h^n)$，意为充分小时 $|R|\le C|h|^n$，并不要求 $R/h^n$ 有非零极限。记 $R=o(h^n)$，则要求该比值趋于零。物理近似通常先确定无量纲小参数，例如 $v/c$、$\ell/r$、摆角 $\theta$，再决定保留到哪一阶。

泰勒系数来自逐阶匹配。先设 $f(x_0+h)\simeq a_0+a_1h+a_2h^2$，在 $h=0$ 处分别匹配函数值、一阶导数和二阶导数，就得到 $a_0=f(x_0)$、$a_1=f'(x_0)$、$2a_2=f''(x_0)$。阶乘正是多项式反复求导时产生的系数。这样得到的多项式优先匹配展开点附近的局部性质，不保证在远处仍有相同形状。
\fig{c3k-remainder}{余项是同一自变量处函数与近似式的差}{\begin{tikzpicture}[x=1.5cm,y=.8cm]
\draw[axis](-.5,0)--(2.1,0)node[right]{$h$};
\draw[axis](0,0)--(0,4.3)node[above]{$f$};
\draw[world,domain=-.35:1.3,samples=40]plot(\x,{exp(\x)});
\draw[blue,domain=-.35:1.6]plot(\x,{1+\x});
\draw[guide](1.2,0)--(1.2,{exp(1.2)});
\draw[dim,red](1.65,2.2)--(1.65,{exp(1.2)})node[midway,right]{$R_1$};
\draw[guide](1.2,2.2)--(1.65,2.2);\draw[guide](1.2,{exp(1.2)})--(1.65,{exp(1.2)});
\fill(0,1)circle(1.5pt)node[left]{$1$};
\node[above left]at(1.1,3){$\e^h$};\node[below,blue]at(1.1,2.1){$1+h$};
\end{tikzpicture}}
图 \ref{c3k-remainder} 中函数与一次近似在展开点有相同高度和斜率，但离开该点后逐渐分开。若要保证误差不超过某一数值，应在整个考察区间内给出高阶导数的上界，再代入余项公式，不能只用展开点处的高阶导数代替这个上界。

含量纲变量时，常先写成 $f(x_0+\ell\epsilon)$，让 $\epsilon$ 成为无量纲小参数。例如长度 $r$ 受微小改变量 $\eta$ 影响，应展开 $\eta/r$，而不是没有比较尺度地声称“$\eta$ 很小”。比较尺度一变，可忽略的项也可能改变。

\subsection{从导数得到展开，并确定有效范围}
以零点为展开中心时，各项系数由函数在零点的逐阶导数确定。指数函数每阶导数仍为 $\e^x$，在零点都等于一，因而
\[
 \e^x=\sum_{n=0}^{\infty}\frac{x^n}{n!}
 =1+x+\frac{x^2}{2!}+\frac{x^3}{3!}+\cdots.
\]
若只保留到三次项，余项满足 $|R_3|\le\e^{|x|}|x|^4/4!$。因此“级数对全部实数收敛”与“少量几项已经足够准确”是两件不同的事：当 $|x|$ 很大时，前几项通常不能给出好的近似。

正弦函数的导数依次为 $\cos x,-\sin x,-\cos x,\sin x$，每四次重复。代入零点，偶数阶导数为零，奇数阶导数交替为 $1,-1$，于是
\[
 \sin x=\sum_{n=0}^{\infty}\frac{(-1)^nx^{2n+1}}{(2n+1)!}
 =x-\frac{x^3}{3!}+\frac{x^5}{5!}-\cdots.
\]
余弦是偶函数，只有偶次项；其零点导数依次给出
\[
 \cos x=\sum_{n=0}^{\infty}\frac{(-1)^nx^{2n}}{(2n)!}
 =1-\frac{x^2}{2!}+\frac{x^4}{4!}-\cdots.
\]
奇偶性不仅节省计算，还能检查展开：若在余弦的零点展开中出现一次项，必有计算或变量选择错误。正弦和余弦的无穷级数对全部实数收敛，作小角近似时仍应根据第一项被舍去的量估计误差。


图\ref{c3deep:smallangle}显示小角展开中“先看所求量”的含义。对 $\sin\theta$ 本身，直线已经给出首项；若所求的是 $\theta-\sin\theta$，直线近似却会把结果抹掉，必须保留三次项：
\[
 \theta-\sin\theta=\frac{\theta^3}{6}+O(\theta^5).
\]
右图单独画差值，放大后才看清这个三阶量。误差的比较对象应随问题改变：对于 $\sin\theta$，舍去三次项造成相对二阶误差；对于差值，三次项就是主项。
\fig{c3deep:smallangle}{原函数很接近直线，不代表二者的差可以忽略}{\begin{tikzpicture}
\begin{scope}[x=2.1cm,y=2cm]\draw[axis](0,0)--(1.35,0)node[right]{$\theta$};\draw[axis](0,0)--(0,1.5)node[above]{$f$};
\draw[blue,thick,domain=0:1.2,samples=70]plot(\x,{sin(\x r)});\draw[red](0,0)--(1.2,1.2);\node[red,right]at(1.2,1.2){$\theta$};\node[blue,right]at(1.2,.932){$\sin\theta$};\end{scope}
\begin{scope}[xshift=5.5cm,x=2.1cm,y=8cm]\draw[axis](0,0)--(1.35,0)node[right]{$\theta$};\draw[axis](0,0)--(0,.38)node[above]{差值};
\draw[blue,thick,domain=0:1.2,samples=70]plot(\x,{\x-sin(\x r)});\draw[red,dashed,domain=0:1.2,samples=70]plot(\x,{(\x)^3/6});
\node[above]at(.85,.3){$\theta-\sin\theta\simeq\theta^3/6$};\end{scope}
\end{tikzpicture}}


对 $f(x)=(1+x)^\alpha$，连续求导给出
$f'(x)=\alpha(1+x)^{\alpha-1}$、
$f''(x)=\alpha(\alpha-1)(1+x)^{\alpha-2}$。
把零点函数值与这两个导数值代入泰勒公式，便得到
\[
 (1+x)^\alpha
 =1+\alpha x+\frac{\alpha(\alpha-1)}2x^2+O(x^3).
\]
逐阶求导的系数遵循同一规律。定义 $\binom\alpha0=1$、$\binom\alpha n=\alpha(\alpha-1)\cdots(\alpha-n+1)/n!$，便可写成
\[
 (1+x)^\alpha=\sum_{n=0}^{\infty}\binom\alpha n x^n,\qquad |x|<1.
\]
若 $\alpha$为非负整数，系数在有限项后全部为零，所得多项式恒等式对任意实数 $x$成立。实数非整数幂在此取 $1+x>0$ 的分支，其无穷展开的端点须另作收敛判断。相对论动能、远处电场与几何长度的低阶近似通常要求 $|x|\ll1$；先化成上述形式，才能明确小参数和保留阶数。

几何级数还能提供一个带显式余项的例子。有限求和恒等式
\[
 \frac1{1-x}=\sum_{j=0}^{N}x^j+\frac{x^{N+1}}{1-x}
\]
对 $x\ne1$ 成立。只有在 $|x|<1$ 时，最后的余项才随 $N\to\infty$ 趋于零，得到无穷展开。将变量换成$-t$，先保留有限和，再从$0$积到$x$：
\[
\int_0^x\frac{\dd t}{1+t}
=\sum_{j=0}^{N}(-1)^j\int_0^xt^j\,\dd t
+\int_0^x\frac{(-t)^{N+1}}{1+t}\,\dd t.
\]
左侧原函数为$\ln(1+t)$，右侧每个多项式项的原函数为$t^{j+1}/(j+1)$。记最后一个积分为$R_N(x)$，则
\begin{align*}
\ln(1+x)
&=[\ln(1+t)]_0^x\\
&=\sum_{j=0}^{N}(-1)^j
\left[\frac{t^{j+1}}{j+1}\right]_0^x+R_N(x)\\
&=\sum_{j=0}^{N}\frac{(-1)^jx^{j+1}}{j+1}+R_N(x).
\end{align*}
在任意闭区间$|x|\le q<1$内，$1+t\ge1-q>0$，因此
\[
|R_N(x)|\le\frac1{1-q}\int_0^{|x|}s^{N+1}\,\dd s
=\frac{|x|^{N+2}}{(1-q)(N+2)}\longrightarrow0.
\]
取$N\to\infty$，才得到$\ln(1+x)=x-x^2/2+x^3/3-\cdots$。这说明逐项积分的依据是余项的一致控制，不能把形式上的无穷求和任意移过积分号。

决定保留阶数时，应观察最后要求的量，而不是只看中间式。例如单摆的切向恢复力第一项与 $\theta$ 成正比，保留一阶就能得到线性方程；其势能 $mgl(1-\cos\theta)$ 的常数项抵消，最低非零项为二阶，因此只保留势能的一阶会把恢复作用整个丢掉。两个相近场强相减也会消去零阶，必须计算到第一个不相消的阶数。

对相近量相减的表达式，先变形再展开有时更稳妥。例如
$1-\sqrt{1-\epsilon^2}=\epsilon^2/[1+\sqrt{1-\epsilon^2}]$，立即可见最低非零阶为 $\epsilon^2$。若分别把两项都粗略取为一，便会误判结果为零。还应区分绝对误差与相对误差：原量本身很小时，一个绝对很小的舍弃项仍可能占据很大比例。

\subsection{误差要求与截断阶数}
实际计算很少需要无穷多项。若要用 $\sin x\simeq x$且要求相对误差小于 $\eta$，对小的非零 $x$，第一项误差的相对量级为 $x^2/6$，所以可先以 $|x|\lesssim\sqrt{6\eta}$估计允许角度，再用严格余项界核验。角度必须使用弧度；将度数直接代入会使小参数失去正确尺度。

图\ref{c3k-order}比较 $\cos x$与其二阶、四阶近似。三条曲线在零点附近贴近，并不意味着离开零点后仍可互换。若所求是 $1-\cos x$，常数一已经抵消，二阶项成为主项；若又从中减去 $x^2/2$，便必须计算到四阶。这种先检查抵消再定阶的次序，比机械地对所有中间式都保留一阶可靠。
\fig{c3k-order}{零点附近的余弦与不同阶数近似}{\begin{tikzpicture}[x=1.8cm,y=1.4cm]
\draw[axis](-1.8,0)--(1.9,0)node[right]{$x$};\draw[axis](0,-.6)--(0,1.5)node[above]{$f$};
\draw[world]plot[domain=-1.6:1.6,samples=80]({\x},{cos(\x r)});
\draw[blue]plot[domain=-1.6:1.6,samples=60]({\x},{1-(\x)^2/2});
\draw[red,dashed]plot[domain=-1.6:1.6,samples=60]({\x},{1-(\x)^2/2+(\x)^4/24});
\node at(0,-.9){黑：$\cos x$\quad 蓝：二阶\quad 红：四阶};
\end{tikzpicture}}

\subsection{先无量纲化，再比较主导项}
设精确量为 $Q(\varepsilon)=Q_0+Q_1\varepsilon+Q_2\varepsilon^2+\cdots$，其中 $\varepsilon$必须无量纲。若 $Q_0\ne0$，省略一阶项的相对误差通常为 $O(\varepsilon)$；若 $Q_0=0$而 $Q_1\ne0$，同样的绝对二阶误差相对主项仍为 $O(\varepsilon)$。因此“误差是二阶”必须说明比较的是绝对量还是相对量。

对两相近力相减的偶极模型，单个场强的主项为 $O(r^{-2})$，但两个主项方向相反，合场的最低非零阶为 $O(d/r^3)$。若每个场只算到零阶，相减会得到零，物理上却丢掉了真正的领先项。反过来，若零阶不抵消，保留不必要的许多高阶只会掩盖主导关系。

一个有效的计算顺序是先把长度写成 $r\sqrt{1+\delta}$，辨明 $\delta$中一次和二次小量；再展开外层函数，最后将各项相加或相减。对 $1/\sqrt{1+\delta}$，二阶系数为 $3/8$，而对 $1/(1+\delta)$，二阶系数为 $1$；二者对应不同外层函数，不能只记“分母展开”而混用系数。




图\ref{c3deep:cancel}给出相消情形的具体几何模型。在轴线上远处的观察点 $P$，两电荷的场都近似与 $q/r^2$ 同阶，但方向相反。保留两距离的区别，才得到
\begin{align*}
 E_x&=\mathcal Kq\left[\frac1{(r-d/2)^2}-\frac1{(r+d/2)^2}\right]\\
 &=\frac{\mathcal Kq}{r^2}\left[\left(1+\frac dr+O((d/r)^2)\right)-\left(1-\frac dr+O((d/r)^2)\right)\right]\\
 &\simeq\frac{2\mathcal Kqd}{r^3},\qquad d/r\ll1.
\end{align*}
这里 $\mathcal K$ 是库仑常量，所需电学输入只有点电荷场随距离平方反比且按方向叠加。两项的零阶部分抵消，首个不抵消项由电荷间距决定；这个结论来自图中两距离的差异，不能仅凭“远处”把两距离先设成完全相等。
\fig{c3deep:cancel}{先保留两条距离的差，再叠加方向相反的贡献}{\begin{tikzpicture}
\draw[axis](-.7,0)--(8,0)node[right]{$x$};\fill[blue](0,0)circle(2.5pt)node[below]{$-q$};\fill[red](1.5,0)circle(2.5pt)node[below]{$+q$};\fill(6.5,0)circle(2pt)node[below]{$P$};
\draw[dim](0,-.65)--(1.5,-.65)node[midway,below]{$d$};\draw[dim](.75,-1.3)--(6.5,-1.3)node[midway,below]{$r$};
\draw[force,red](6.5,.4)--(7.8,.4)node[above]{$E_+$};\draw[force,blue](6.5,1)--(5.5,1)node[left]{$E_-$};\node at(4,1.8){$|E_+|>|E_-|$};
\end{tikzpicture}}


\subsection{复指数与振动表示}
令 $\mathrm i^2=-1$。将绝对收敛的指数级数分成偶数项与奇数项，
\[
 \e^{\mathrm i\theta}
 =\left(1-\frac{\theta^2}{2!}+\cdots\right)
 +\mathrm i\left(\theta-\frac{\theta^3}{3!}+\cdots\right)
 =\cos\theta+\mathrm i\sin\theta.
\]
因此正弦振动可以由复指数的实部或虚部表示。复表示是一种运算工具，实际位移仍取实数部分；振幅、相位和频率应在取实部之前说明。遇到阻尼振动时，指数实部给出衰减，虚部给出振荡。

例如 $x(t)=A\cos(\omega t+\phi)$ 可以写成
$x(t)=\operatorname{Re}[A\e^{\mathrm i\phi}\e^{\mathrm i\omega t}]$。对复指数求一次时间导数相当于乘以 $\mathrm i\omega$，求两次则乘以 $-\omega^2$；取实部后恰好得到 $\ddot x=-\omega^2x$。复数的乘法同时记录振幅与相位，因而适合处理线性振动中的叠加。若方程含有非线性项，不能未经核验就把求实部与非线性运算交换。

\fig{c3k-complex}{复平面中的转动对应正弦振动的投影}{\begin{tikzpicture}[scale=.9]
\draw[axis](-.4,0)--(3,0)node[right]{$\operatorname{Re}z$};
\draw[axis](0,-.4)--(0,2.5)node[above]{$\operatorname{Im}z$};
\coordinate(P)at(1.8,1.2);\draw[flow,blue](0,0)--(P)node[midway,above left]{$A$};
\draw[guide](P)--(1.8,0);\draw[guide](P)--(0,1.2);
\draw(.6,0)arc(0:33.69:.6);\node at(.85,.24){$\theta$};
\fill(P)circle(1.6pt)node[above right]{$z=A\e^{\mathrm i\theta}$};
\node[below]at(1.8,0){$A\cos\theta$};\node[left]at(0,1.2){$A\sin\theta$};
\end{tikzpicture}}
图 \ref{c3k-complex} 中，模长 $A$ 决定投影的振幅，辐角 $\theta=\omega t+\phi$ 决定投影的相位。乘以 $\mathrm i$ 相当于在复平面旋转直角，所以求导不仅带来系数 $\omega$，还使相位超前 $\pi/2$；二次求导再转直角，便与原位移反向。

处理多个同频线性振动时，可以先把各自的复振幅 $A_j\e^{\mathrm i\phi_j}$ 相加，求合成复振幅的模与辐角，最后取实部。不同频率的项不能合并成一个恒定复振幅；非线性作用更不能先随意合并再取实部。复数表示简化了运算，但没有改变实际运动必须满足实数方程与初值的要求。

\section{曲线切线、曲率与运动}
\subsection{曲率测量方向改变的快慢}
切线方向角为 $\psi$，弧长为 $s$。曲率定义为 $\kappa=|\dd\psi/\dd s|$，曲率半径 $\rho=1/\kappa$。直线的方向不变，曲率为零；圆在任何点都有 $\kappa=1/R$。一般曲线在一点附近可以用具有相同切线、相同曲率的密切圆作二阶近似。

若曲线写成 $y=f(x)$，则 $\tan\psi=y'$。对 $x$ 求导，先得到
$(1+\tan^2\psi)\,\dd\psi/\dd x=y''$，即
$\dd\psi/\dd x=y''/(1+y'^2)$。
另一方面，小段路径的水平、竖直分量分别为 $\dd x,y'\dd x$，由勾股关系有 $\dd s=\sqrt{1+y'^2}\,|\dd x|$。把转角变化除以弧长变化，才得到
\[
 \kappa=\frac{|y''|}{(1+y'^2)^{3/2}},\qquad
 \boxed{\rho=\frac{(1+y'^2)^{3/2}}{|y''|}}.
\]
这里撇号是对 $x$ 求导。$y''=0$ 时不能机械地说存在一个有限曲率半径；在普通光滑情形曲率为零、半径趋于无穷。

\fig{c3k-curvature}{沿弧长行进时切线转过的角度}{\begin{tikzpicture}[scale=.95]
\coordinate(O)at(0,0);\coordinate(P)at(2.125,.569);\coordinate(Q)at(1.1,1.905);
\draw[world](2.2,0)arc(0:85:2.2);
\draw[very thick,blue](P)arc(15:60:2.2);
\draw[guide](O)--(P);\draw[guide](O)--(Q);
\draw[flow](P)--++(-.285,1.063)node[right]{$\bm e_\tau(P)$};
\draw[flow](Q)--++(-.953,.55)node[above]{$\bm e_\tau(Q)$};
\draw(.58,.155)arc(15:60:.6);\node at(.88,.6){$\Delta\psi$};
\fill(P)circle(1.5pt)node[below right]{$P$};\fill(Q)circle(1.5pt)node[above right]{$Q$};
\node[below]at(O){$O$};\node[right,blue]at(2.9,1.05){$\Delta s$};
\draw[guide,blue](2.9,1.05)--(1.86,1.18);
\end{tikzpicture}}
图 \ref{c3k-curvature} 用圆弧说明两个切向量的方向差。切线均与半径垂直，因此切线转角等于半径转角，$\Delta s=R\Delta\psi$，给出 $\kappa=1/R$。一般曲线则先取很短的一段，再求转角与弧长的比值极限。曲率大表示同样弧长内转向更明显，与图画得大还是小应作物理尺度上的区分。

曲率公式中的 $x,y$ 应是同一空间中的长度坐标。把位置对时间的图线直接当成实际空间轨道，会混淆坐标单位：图线的凹凸能表示加速度符号，却不能直接解释为质点轨道的曲率。对实际轨迹，改变绘图时横纵轴的比例会改变纸面形状，计算仍必须使用实际坐标与一致的长度单位。

\subsection{密切圆的位置也由局部导数确定}
在平面正则点，曲率半径只给出密切圆的大小，圆心还必须沿曲线内法线定位。若轨迹 $y=f(x)$且 $f''\ne0$，取向上法线单位矢量 $(-f',1)/\sqrt{1+f'^2}$。有符号曲率为 $f''/(1+f'^2)^{3/2}$，故圆心坐标为
\[x_C=x-\frac{f'(1+f'^2)}{f''},\qquad
y_C=f(x)+\frac{1+f'^2}{f''}.\]
$f''>0$时圆心在向上法线侧，$f''<0$时在另一侧。曲率公式中的绝对值只用于半径大小，若在圆心公式中也随意取绝对值，就会把弯曲方向画反。

对抛物线 $y=ax^2$的顶点，$f'=0$、$f''=2a$，所以圆心为 $(0,1/(2a))$，半径为 $1/(2a)$。这里抛物线焦点高度为 $1/(4a)$，与密切圆心高度不同；焦点控制全局反射性质，曲率中心描述局部二阶弯曲，不能把两种几何点混同。



图\ref{c3deep:osculating}把抛物线顶点 $O$、焦点 $F$ 和密切圆心 $C$ 分开标出。圆在 $O$ 的下半支可写成
\[
 y=\rho-\sqrt{\rho^2-x^2}=\frac{x^2}{2\rho}+O(x^4/\rho^3).
\]
只有取 $\rho=1/(2a)$，二次项才与 $y=ax^2$ 相同。这既验证密切圆半径，也说明密切圆只在局部近似曲线；它不会代替整个抛物线的反射性质。
\fig{c3deep:osculating}{焦点与局部曲率中心是不同的点}{\begin{tikzpicture}[scale=1.2]
\draw[axis](-2.1,0)--(2.1,0)node[right]{$x$};\draw[axis](0,-.25)--(0,3.5)node[above]{$y$};
\draw[red,dashed](0,1.5)circle(1.5);\draw[blue,thick,domain=-2:2,samples=80]plot(\x,{(\x)^2/3});
\fill(0,0)circle(1.6pt)node[below left]{$O$};\fill(0,.75)circle(1.6pt)node[right]{$F:\ 1/(4a)$};\fill(0,1.5)circle(1.6pt)node[right]{$C:\ 1/(2a)$};
\draw[dim](-.45,0)--(-.45,1.5)node[midway,left]{$\rho$};\node[blue,right]at(1.7,.96){$y=ax^2$};
\end{tikzpicture}}


\subsection{用运动学计算曲率}
任意正则参数 $t$ 下，
\[
 \rho=\frac{(\dot x^2+\dot y^2)^{3/2}}
 {|\dot x\ddot y-\dot y\ddot x|}.
\]
推导时，平面内的叉积大小为
$|\bm v\times\bm a|=|\dot x\ddot y-\dot y\ddot x|$。
切向加速度与速度平行，不贡献叉积，故该大小等于 $v a_n$。先除以 $v$ 得 $a_n=|\bm v\times\bm a|/v$，再代入 $\rho=v^2/a_n$，就得到 $\rho=v^3/|\bm v\times\bm a|$，也就是上式。曲率由路径决定，与沿曲线快走、慢走无关；若参数在某点使速度为零，应换用正则参数或取极限。

特别地，椭圆参数 $x=a\cos\theta,y=b\sin\theta$ 便于求导，但若 $\dot\theta$ 恒定，实际速率一般并不恒定。只有在特定顶点，或另行证明切向加速度为零时，才可把全部加速度直接代入 $v^2/\rho$。

这个表达式对正则参数的选择不敏感。若改用参数 $q$，则
\[
 \dot{\bm r}=\bm r_q\dot q,\qquad
 \ddot{\bm r}=\bm r_{qq}\dot q^2+\bm r_q\ddot q.
\]
叉积中含 $\bm r_q\times\bm r_q$ 的项为零，所以
$|\dot{\bm r}\times\ddot{\bm r}|=|\bm r_q\times\bm r_{qq}|\,|\dot q|^3$；
分子的 $v^3=|\bm r_q|^3|\dot q|^3$ 带有相同因子，约去后得到只依赖曲线的结果。参数匀速并不是推导条件，参数在该点能够正常描绘曲线才是条件。

因此利用运动求曲率时，应先分清真正的切向与法向加速度。若只知道合加速度大小，可以用 $a_n=\sqrt{a^2-\dot v^2}$，或直接用叉积排除切向部分。图 \ref{c3k-curvature} 中的转向只由法向变化决定；沿同一条路径加快行进虽改变加速度大小，却不会改变该位置的密切圆。

% thirty-c3-theory

\section{多变量的变化与空间方向}
\subsection{偏导数保持其余自变量不变}
一个量有时同时依赖多个独立变量，例如场点的电势取决于位置，气体压强取决于温度和体积。对可微函数 $f(x,y)$，询问“变化率”之前，必须说明沿哪个方向改变自变量。固定 $y$，只改变 $x$，得到的偏导数定义为
\[
\frac{\partial f}{\partial x}(x,y)
=\lim_{h\to0}\frac{f(x+h,y)-f(x,y)}h.
\]
同理，$\partial f/\partial y$ 的计算固定 $x$。符号 $\partial$ 用来强调这个约束。比如 $f=x^2y$，对 $x$ 求偏导时 $y$ 是常量，得 $2xy$；对 $y$ 求偏导时 $x^2$ 是常量，得 $x^2$。

图\ref{c3k:partial}把从 $P$ 到邻近点 $Q$ 的变化分成两段：$P\to A$ 只改变 $x$，$A\to Q$ 只改变 $y$。将两段函数增量相加，在导数连续的条件下保留一阶项，得到
\[
\Delta f=f_x\Delta x+f_y\Delta y
 +o\!\left(\sqrt{\Delta x^2+\Delta y^2}\right),\qquad
\dd f=f_x\dd x+f_y\dd y.
\]
这里 $f_x,f_y$ 均在起点计算。两条微小路径次序不同引起的差别属于更高阶量；全微分给的是各方向一阶变化的合成。
\fig{c3k:partial}{沿两个独立坐标方向分解函数的增量}{\begin{tikzpicture}
\draw[axis](-.3,0)--(6,0)node[right]{$x$};\draw[axis](0,-.3)--(0,3.8)node[above]{$y$};
\coordinate(P)at(1,1);\coordinate(A)at(4.5,1);\coordinate(Q)at(4.5,3);
\draw[guide](1,0)--(P)--(0,1);\draw[guide](4.5,0)--(A);\draw[guide](Q)--(0,3);
\draw[flow,blue](P)--(A)node[midway,below]{$\dd x$，$y$ 不变};
\draw[flow,red](A)--(Q)node[midway,right]{$\dd y$，$x$ 不变};
\draw[flow](P)--(Q)node[midway,above left]{$\dd\bm r$};
\fill(P)circle(1.5pt)node[left]{$P$};\fill(A)circle(1.5pt)node[below right]{$A$};\fill(Q)circle(1.5pt)node[above]{$Q$};
\end{tikzpicture}}

若 $x,y$ 又沿某条运动路径随 $t$ 变化，把 $\dd x=\dot x\dd t$、$\dd y=\dot y\dd t$ 代入，得到
\[
\frac{\dd f}{\dd t}=\frac{\partial f}{\partial x}\dot x
 +\frac{\partial f}{\partial y}\dot y.
\]
偏导描述单独改变一个自变量的效应，沿路径的总变化率还必须计入两个变量实际怎样变化。只算其中一项，会漏掉另一段位移的贡献。

\subsection{方向导数要除以真实的位移长度}
在直角坐标中，$\dd\bm r=\dd x\bm e_x+\dd y\bm e_y$，所以可将全微分写成
\[
\dd f=\bm\nabla f\cdot\dd\bm r,\qquad
\bm\nabla f=f_x\bm e_x+f_y\bm e_y.
\]
向量 $\bm\nabla f$ 称为梯度。沿单位方向 $\bm n$ 走过小距离 $\dd s$，有 $\dd f/\dd s=\bm\nabla f\cdot\bm n$；它把某个方向上的实际距离变化率从全微分中取出。

换成半径 $r$ 和角度 $\theta$ 时，角度本身不是长度。图\ref{c3k:polar-increment}在一个子午面内从 $+z$ 轴量 $\theta$：$P\to P_r$ 固定角度，其长度为 $\dd r$；$P\to P_\theta$ 固定半径，其弧长为 $r\dd\theta$。角度用弧度，两种位移在一阶上正交。
\fig{c3k:polar-increment}{径向线段与增角圆弧对应不同的长度因子}{\begin{tikzpicture}
\draw[axis](0,-.5)--(0,3.7)node[above]{$z$};\draw[axis](-.3,0)--(4.6,0)node[right]{$x$};
\coordinate(P)at(30:2.6);\coordinate(R)at(30:3.5);\coordinate(T)at(5:2.6);
\draw[guide](0,0)--(30:3.9);\draw[guide](0,0)--(T);
\draw[world](0,0)--(P)node[midway,above left]{$r$};
\draw[flow,blue](P)--(R)node[midway,above]{$\dd r$};
\draw[flow,red](P)arc(30:5:2.6)node[midway,right,xshift=6pt]{$r\dd\theta$};
\draw(0,.65)arc(90:30:.65);\node at(60:.95){$\theta$};
\draw[flow](30:1.25)arc(30:5:1.25);\node at(17.5:1.7){$\dd\theta$};
\fill(P)circle(1.4pt)node[above left]{$P$};\fill(R)circle(1.4pt)node[right]{$P_r$};\fill(T)circle(1.4pt)node[below right]{$P_\theta$};\node[below left]at(0,0){$O$};
\end{tikzpicture}}
因此，在 $r>0$ 的这两个方向上，
\[
\dd f=f_r\dd r+f_\theta\dd\theta,
\qquad
\bm\nabla f=f_r\bm e_r+\frac1r f_\theta\bm e_\theta.
\]
这里的 $1/r$ 来自把“每单位角度的变化”换成“每单位弧长的变化”。例如静电场满足 $\dd V=-\bm E\cdot\dd\bm r$，便有 $E_r=-V_r$、$E_\theta=-V_\theta/r$。若三维问题还随方位角变化，就必须加入第三个方向；轴对称时才可只保留图中的两个方向。在坐标原点应改用不奇异的直角坐标或适当极限。

\section{例题与分析}
\begin{question}
\begin{example}[驻点与极值的分类]
在各自实数定义域内，分析 $x^2$、$x^2+1/x^2$、$x-2\sin x$、$x^3$、$\sqrt{1-x^2}$ 的极值和全局最值。
\end{example}
\fig{c3e1q}{导数为零不一定形成峰或谷}{\begin{tikzpicture}[x=1cm,y=.7cm]
\draw[axis](-1.5,0)--(1.6,0);\draw[axis](0,-.2)--(0,2.2);\draw[blue,domain=-1.2:1.2,samples=40]plot(\x,{\x*\x});\node[above]at(0,2.3){$x^2$};
\begin{scope}[xshift=4cm]\draw[axis](-1.5,0)--(1.6,0);\draw[axis](0,-1.6)--(0,1.8);\draw[blue,domain=-1.1:1.1,samples=40]plot(\x,{(\x)^3});\node[above]at(0,2.3){$x^3$};\end{scope}
\end{tikzpicture}}
\end{question}
\begin{solution}\textbf{比较驻点两侧的导数符号}\par\nobreak
图\ref{c3e1q}中，对$x^2$求导得到$2x$：$x<0$时为负，$x>0$时为正，故零点由减转增，有全局最小值零；$|x|\to\infty$时函数无界，故无最大值。对$x^3$求导得到$3x^2\ge0$，零点两侧不变号，故无局部极值；函数在两端分别趋于正、负无穷，也无全局最值。

对$x^2+x^{-2}$，定义域排除零。先按幂函数法则逐项求导，
\[
\frac{\dd}{\dd x}(x^2+x^{-2})=2x-2x^{-3},\qquad
\frac{\dd}{\dd x}(2x-2x^{-3})=2+6x^{-4}>0.
\]
令一阶导数为零，$2x-2/x^3=0$两边乘以非零的$x^3/2$，得到$x^4=1$，实根只有$\pm1$。二阶导数均为正，两点均为极小；代回原函数，值均为$1+1=2$。全局判断再用
$x^2+x^{-2}-2=(x-1/x)^2\ge0$，证明二为全局下界。$|x|\to\infty$ 时平方项无界，故没有最大值。
\fig{c3e1s}{在 $[-\pi,\pi]$ 内用导数符号判定极值}{\begin{tikzpicture}
\draw[axis](0,0)--(6.7,0)node[right]{$x$};
\foreach \x in {.4,2.3,4.2,6.1}{\draw(\x,-.06)--(\x,.06);}
\node[below]at(.4,0){$-\pi$};\node[below]at(2.3,0){$-\pi/3$};\node[below]at(4.2,0){$\pi/3$};\node[below]at(6.1,0){$\pi$};
\node[above]at(1.35,.2){$f'>0$};\node[above]at(3.25,.2){$f'<0$};\node[above]at(5.15,.2){$f'>0$};
\node[below=17pt]at(2.3,0){极大};\node[below=17pt]at(4.2,0){极小};
\end{tikzpicture}}



\par\medskip\Needspace{4\baselineskip}\noindent\textbf{区分周期性局部极值与全局最值}\par\nobreak
对 $f=x-2\sin x$，$f'=1-2\cos x$，驻点要求 $\cos x=1/2$，所以 $x=2n\pi\pm\pi/3$。二阶导数为 $2\sin x$：在减号支等于 $-\sqrt3<0$，在加号支等于 $\sqrt3>0$。因此图 \ref{c3e1s} 中，减号支为极大点，其值为
$(2n\pi-\pi/3)-2(-\sqrt3/2)=2n\pi-\pi/3+\sqrt3$；加号支为极小点，其值为 $2n\pi+\pi/3-\sqrt3$。由于 $x-2\le f(x)\le x+2$，函数随 $x\to\pm\infty$ 分别无界地增大、减小，没有全局最值。

最后一式定义域是 $[-1,1]$。内部按链式法则求导：
\[
f'(x)=\frac12(1-x^2)^{-1/2}(-2x)
=-\frac{x}{\sqrt{1-x^2}}.
\]
导数在 $(-1,0)$ 为正、$(0,1)$ 为负，因此零点为极大点。比较实际函数值：$f(0)=1$，$f(\pm1)=0$，又因整个区间内 $0\le\sqrt{1-x^2}\le1$，所以最大值为一、最小值为零。端点不要求导数为零，仍须参加全局最值比较。
\end{solution}

\begin{question}
\begin{example}[圆锥上的弹性绳套]
质量 $m$、自然长度 $L$、整体劲度系数 $k$ 的均匀弹性绳套水平套在光滑圆锥上，圆锥顶角为 $2\alpha$。绳套保持水平圆形，只作升降。求距顶点的平衡高度 $z$，并讨论该方向稳定性。
\end{example}
\fig{c3e2q}{向下坐标决定绳套的周长}{\begin{tikzpicture}
\draw[world](-1.7,0)--(0,2.8)--(1.7,0);\draw[axis](0,3.1)--(0,-.3)node[below]{$z$};
\draw[blue,very thick](0,1.1)ellipse(1.032 and .22);\draw[guide](0,1.1)--(1.032,1.1);
\node[right]at(1.2,1.1){$r=z\tan\alpha$};
\draw(0,2.1)arc(-90:-58.74:.7);\node at(.23,1.83){$\alpha$};
\end{tikzpicture}}
\end{question}
\begin{solution}\textbf{用下降高度统一绳长与势能}\par\nobreak
图 \ref{c3e2q} 中向下为正，圆环半径 $r=z\tan\alpha$，伸长量为 $2\pi z\tan\alpha-L$。重力势能为 $-mgz$，故
\[
 U(z)=\frac k2(2\pi z\tan\alpha-L)^2-mgz.
\]
对$z$求导时，$k,L,\alpha,m,g$均为固定参数。先按平方的链式法则写出
\[
U'(z)=\frac k2\,2(2\pi z\tan\alpha-L)(2\pi\tan\alpha)-mg.
\]
令$U'(z_0)=0$，整理后得到
\[
 2\pi k\tan\alpha(2\pi z_0\tan\alpha-L)=mg,
\]
先将平衡方程除以正数 $2\pi k\tan\alpha$，得到
$2\pi z_0\tan\alpha-L=mg/(2\pi k\tan\alpha)$，再移项并除以 $2\pi\tan\alpha$，就得到平衡高度。
\[
 \boxed{z_0=\frac{L}{2\pi\tan\alpha}
 +\frac{mg}{4\pi^2k\tan^2\alpha}}.
\]
\fig{c3e2s}{平衡位置附近形成势阱}{\begin{tikzpicture}
\draw[axis](0,0)--(5,0)node[right]{$z$};\draw[axis](0,0)--(0,2.7)node[above]{$U$};
\draw[blue,domain=.5:4.5,samples=40]plot(\x,{.45*(\x-2.5)^2+.4});\draw[guide](2.5,0)--(2.5,.4);\node[below]at(2.5,0){$z_0$};
\draw[flow](3.6,1.3)--(2.9,.8);
\end{tikzpicture}}



\par\medskip\Needspace{4\baselineskip}\noindent\textbf{检查张紧条件与恢复作用}\par\nobreak
如图\ref{c3e2s}，再对$U'(z)$中的线性项求导，
\[
U''(z)=2\pi k\tan\alpha\,(2\pi\tan\alpha)-0
=4\pi^2k\tan^2\alpha>0.
\]
所以升降扰动后存在恢复作用。平衡伸长量为$mg/(2\pi k\tan\alpha)>0$，也验证绳套处于张紧状态。若圆锥有限长，还须检查$z_0$未越过底边。


\par\medskip\Needspace{4\baselineskip}\noindent\textbf{由绳元的实际速度确定振动惯量}\par\nobreak
% thirty-c3-conekinetic
若进一步研究保持水平圆形的升降振动，不能直接把 $z$的有效质量写成 $m$。因为绳套下降时半径也增加，$r=z\tan\alpha$给 $\dot r=\dot z\tan\alpha$。每个绳元同时具有竖直和径向速度，图\ref{c3e2velocity}给出
\[v_{\text{元}}^2=\dot z^2+\dot r^2=\dot z^2\sec^2\alpha,
\qquad T=\frac12m\sec^2\alpha\dot z^2.\]
\fig{c3e2velocity}{圆锥母线上运动的绳元具有两个速度分量}{\begin{tikzpicture}
\draw[world](0,2.7)--(1.7,0);\coordinate(P)at(.85,1.35);
\draw[flow,blue](P)--++(0,-1)node[below]{$\dot z$};
\draw[flow,red](.85,.35)--(1.47963,.35)node[right]{$\dot r=\dot z\tan\alpha$};
\draw[flow](P)--(1.47963,.35)node[midway,right]{$\bm v_{\text{元}}$};
\fill(P)circle(1.3pt)node[left]{绳元};\node at(.3,2){$\alpha$};
\end{tikzpicture}}
令 $z=z_0+\eta$，势能在平衡位置附近为 $U_0+2\pi^2k\tan^2\alpha\,\eta^2$。因此
\[m\sec^2\alpha\ddot\eta+4\pi^2k\tan^2\alpha\,\eta=0.\]
代入简谐位移的$\ddot\eta=-\omega^2\eta$，非零振动要求
\[
-m\sec^2\alpha\,\omega^2+4\pi^2k\tan^2\alpha=0.
\]
于是$\omega^2=4\pi^2k\tan^2\alpha/(m\sec^2\alpha)
=(4\pi^2k/m)\sin^2\alpha$。
这一步说明，势能的二阶系数确定回复作用，而频率还必须结合实际速度场确定惯量。若只取竖直速度，就会漏掉径向运动的动能。
\end{solution}

\begin{question}
\begin{example}[弹簧连接的引力双球]
两球各质量 $M$，在同一直线上由轻弹簧连接，仅考虑相互引力与弹力。平衡时压缩量为 $x>0$，劲度系数为 $k$，球半径忽略。求弹簧自然长度，并讨论沿连线方向的稳定性。
\end{example}
\fig{c3e3q}{相互吸引与压缩弹簧的排斥}{\begin{tikzpicture}
\fill(0,0)circle(4pt)node[above]{$M$};\fill(4.6,0)circle(4pt)node[above]{$M$};
\draw[world,decorate,decoration={zigzag,segment length=5pt,amplitude=3pt}](.2,0)--(4.4,0);
\draw[dim](0,-1.05)--(4.6,-1.05)node[midway,below]{$r_0$};
\draw[force,red](0,.6)--(.9,.6)node[midway,above]{$F_g$};\draw[force,red](4.6,.6)--(3.7,.6)node[midway,above]{$F_g$};
\draw[force,blue](0,-.35)--(-.9,-.35)node[midway,below]{$kx$};\draw[force,blue](4.6,-.35)--(5.5,-.35)node[midway,below]{$kx$};
\end{tikzpicture}}
\end{question}
\begin{solution}\textbf{由受力平衡确定弹簧自然长度}\par\nobreak
图 \ref{c3e3q} 中平衡时 $kx=GM^2/r_0^2$，因此
\[
 r_0=\sqrt{\frac{GM^2}{kx}},\qquad
 \boxed{L_0=x+\sqrt{\frac{GM^2}{kx}}}.
\]
为判断距离扰动，写系统势能
$U(r)=\tfrac12k(L_0-r)^2-GM^2/r$。
对第一项用链式法则，对第二项写成负幂：$U'(r)=\frac k2\,2(L_0-r)(-1)-GM^2(-r^{-2})=k(r-L_0)+GM^2/r^2$，在 $r=r_0$ 处等于零，恰好恢复受力平衡。再用 $(r^{-2})'=-2r^{-3}$，得到
\[
 U''(r_0)=k-\frac{2GM^2}{r_0^3}
 =k\left(1-\frac{2x}{r_0}\right).
\]
\fig{c3e3s}{距离增加时两种力的变化率相互竞争}{\begin{tikzpicture}
\draw[axis](0,0)--(5.8,0)node[right]{$r$};\node[below]at(2.5,0){$r_0$};
\draw[flow](2.5,.45)--(3.6,.45)node[midway,above]{$\eta>0$};
\node[align=center]at(2.8,1.6){引力减弱率 $2GM^2/r_0^3$\\弹簧排斥力减弱率 $k$};
\end{tikzpicture}}



\par\medskip\Needspace{4\baselineskip}\noindent\textbf{比较两种力的变化率，判定稳定范围}\par\nobreak
图 \ref{c3e3s} 中，若弹簧排斥力减小得更快，拉开后净力会把两球拉回，故稳定条件为
\[
 r_0>2x\quad\Longleftrightarrow\quad
 \boxed{kx^3<GM^2/4}.
\]
把 $r_0=\sqrt{GM^2/(kx)}$ 代入 $r_0>2x$，两侧均正，可以平方，得到 $GM^2/(kx)>4x^2$；再乘 $kx>0$，才得到 $kx^3<GM^2/4$。反向不等号为不稳定。等号处 $U''=0$，但 $U'''=6GM^2/r_0^4>0$，势能首个变化项为 $U'''(r_0)(r-r_0)^3/6$，在平衡点两侧异号，不形成局部极小，仍不稳定。这里 $kx^3$ 与 $GM^2$ 量纲相同。


\par\medskip\Needspace{4\baselineskip}\noindent\textbf{以两球相对坐标建立小振动方程}\par\nobreak
% thirty-c3-gravity
可把稳定结论继续转成小振动方程。两球各自相对质心的位移大小为 $r/2$，速度大小为 $\dot r/2$，所以系统动能为
\[T=2\cdot\frac12M\left(\frac{\dot r}{2}\right)^2
   =\frac12\left(\frac M2\right)\dot r^2.\]
因此距离坐标的有效质量是 $M/2$。在稳定分支令 $r=r_0+\eta$，线性回复项为 $-U''(r_0)\eta$，有
\[\frac M2\ddot\eta+k\left(1-\frac{2x}{r_0}\right)\eta=0.\]
对简谐位移有$\ddot\eta=-\omega^2\eta$；代回并约去非零解的共同因子，得到
\[
-\frac M2\omega^2+k\left(1-\frac{2x}{r_0}\right)=0,
\qquad \omega^2=\frac{2k}{M}\left(1-\frac{2x}{r_0}\right)>0.
\]
图\ref{c3e3s}中当两种力的变化率逐渐接近时，回复系数和频率都减小；在临界等号处，线性频率为零，必须恢复已讨论的三阶势能项，不能据此判为普通的稳定小振动。
\end{solution}

\begin{question}
\begin{example}[圆环上的零自然长度弹簧]
竖直光滑圆环半径 $R$，最高点 $A$ 与质量 $m$ 的小珠 $P$ 之间连接劲度系数 $k$、自然长度为零的轻弹簧。小珠仅能沿圆环运动。求平衡位置及稳定性。
\end{example}
\fig{c3e4q}{用圆心角描述珠的位置}{\begin{tikzpicture}[scale=1.1]
\draw[world](0,0)circle(1.5);\coordinate(A)at(0,1.5);\coordinate(B)at(0,-1.5);\coordinate(P)at(1.299,-.75);
\draw[guide](A)--(B);\draw[guide](0,0)--(P);\draw[world,decorate,decoration={zigzag,amplitude=2pt,segment length=5pt}](A)--(P);
\fill(P)circle(2pt)node[right]{$P$};\node[above]at(A){$A$};\node[below]at(B){$B$};\node[left]at(0,0){$O$};
\draw(0,-.5)arc(-90:-30:.5)node[midway,below right]{$\varphi$};
\end{tikzpicture}}
\end{question}
\begin{solution}\textbf{由弦长与高度写出角坐标势能}\par\nobreak
图 \ref{c3e4q} 中从向下半径起量圆心角 $\varphi$。以圆心为原点，$A=(0,R)$、$P=(R\sin\varphi,-R\cos\varphi)$，所以
\[
 AP^2=R^2\sin^2\varphi+R^2(1+\cos\varphi)^2
 =2R^2(1+\cos\varphi).
\]
最低点高度为 $-R$，珠高于 $B$ 的高度为 $-R\cos\varphi-(-R)=R(1-\cos\varphi)$。因弹簧自然长度为零，其伸长量就是 $AP$，于是
\[
 U=kR^2(1+\cos\varphi)+mgR(1-\cos\varphi)
 =\text{常数}+R(kR-mg)\cos\varphi.
\]
对角坐标求导有
$U'=-R(kR-mg)\sin\varphi$。
当 $kR\ne mg$，平衡要求 $\sin\varphi=0$，在完整圆周上只对应最低点 $B$ 与最高点 $A$（角度相差 $2\pi$ 表示同一位置）。
\fig{c3e4s}{参数变化使势能最低位置交换}{\begin{tikzpicture}[x=.85cm,y=.65cm]
\draw[axis](0,0)--(4,0)node[right]{$\varphi$};\draw[axis](0,0)--(0,3.5)node[above]{$U$};
\draw[blue,domain=0:3.1416,samples=50]plot(\x,{1.8-cos(\x r)});
\draw[red,domain=0:3.1416,samples=50]plot(\x,{1.8+cos(\x r)});
\node[below]at(0,0){$B$};\node[below]at(3.1416,0){$A$};
\node[right,blue]at(3.2,2.8){$kR<mg$};\node[right,red]at(3.2,.8){$kR>mg$};
\end{tikzpicture}}



\par\medskip\Needspace{4\baselineskip}\noindent\textbf{分别检查上下端点及临界参数}\par\nobreak
二阶导数为 $U''=-R(kR-mg)\cos\varphi$。在 $B$，$\cos0=1$，故 $U''_B=R(mg-kR)$；在 $A$，$\cos\pi=-1$，故 $U''_A=R(kR-mg)$。图 \ref{c3e4s} 因而显示：$kR<mg$ 时 $B$ 稳定、$A$ 不稳定；$kR>mg$ 时正好相反。$kR=mg$ 时角度项系数严格为零，势能完全不变，全部位置均为该约束下的随遇平衡。


\par\medskip\Needspace{4\baselineskip}\noindent\textbf{在各稳定点附近求振动频率}\par\nobreak
% thirty-c3-ring
稳定点附近的频率也可由同一势能直接求出。用题设圆心角 $\varphi$，小珠动能为 $T=mR^2\dot\varphi^2/2$。在稳定的最低点 $B$附近令 $\varphi=\eta$，有
\[mR^2\ddot\eta+R(mg-kR)\eta=0.\]
代入$\ddot\eta=-\omega_B^2\eta$，得到$-mR^2\omega_B^2+R(mg-kR)=0$，故
\[\omega_B^2=\frac{R(mg-kR)}{mR^2}
=\frac gR-\frac km\quad(kR<mg).\]
在稳定的最高点$A$附近令$\varphi=\pi+\eta$，此处回复系数为$U''_A=R(kR-mg)$。同样代入简谐位移，
\begin{align*}
-mR^2\omega_A^2+R(kR-mg)&=0,\\
\omega_A^2&=\frac{R(kR-mg)}{mR^2}
=\frac km-\frac gR\quad(kR>mg).
\end{align*}
图\ref{c3e4s}中哪一个点是势阱，便在哪一点使用正的频率平方。临界时 $kR=mg$使整个势能严格为常数，而不只是二阶项消失；此时若给予切向初速，小珠会沿环匀速运动，不会自行返回某个唯一平衡点。这与存在四次势阱但线性频率为零的情形不同。
\end{solution}

\begin{question}
\begin{example}[相对论动能的低速极限]
静质量为 $m_0$ 的粒子总能量为 $E=m_0c^2/\sqrt{1-v^2/c^2}$。求动能的前两项非零展开，并估计经典动能的相对误差首项。
\end{example}
\fig{c3e5q}{小参数来自速度与光速之比}{\begin{tikzpicture}
\draw[axis](0,0)--(6,0)node[right]{$v/c$};\draw[fill=blue!12](0,-.12)rectangle(1.2,.12);
\node[below]at(0,-.15){$0$};\node[below]at(5.3,-.15){$1$};\node[above]at(1.1,.35){低速区};
\end{tikzpicture}}
\end{question}
\begin{solution}\textbf{以速度比的平方作为展开参数}\par\nobreak
图\ref{c3e5q}的低速区满足$\beta=v/c\ll1$，真正进入根式的展开量是$z=\beta^2$。令$f(z)=(1-z)^{-1/2}$，先对$z$逐次求导：
\begin{align*}
f'(z)&=-\frac12(1-z)^{-3/2}(-1)=\frac12(1-z)^{-3/2},\\
f''(z)&=\frac12\left(-\frac32\right)(1-z)^{-5/2}(-1)
=\frac34(1-z)^{-5/2}.
\end{align*}
因此$f(0)=1$、$f'(0)=1/2$、$f''(0)=3/4$，代入二阶泰勒式，
\[
f(z)=1+\frac12z+\frac{3/4}{2!}z^2+O(z^3)
=1+\frac12z+\frac38z^2+O(z^3).
\]
动能必须减去静能，常数项在这一步抵消：
\begin{align*}
K&=m_0c^2\left[1+\frac12\beta^2+\frac38\beta^4
 +O(\beta^6)-1\right]\\
&=\frac12m_0c^2\frac{v^2}{c^2}
 +\frac38m_0c^2\frac{v^4}{c^4}+O(m_0v^6/c^4)\\
&=\boxed{\frac12m_0v^2+\frac38m_0\frac{v^4}{c^2}
 +O(m_0v^6/c^4)}.
\end{align*}
\fig{c3e5s}{经典动能与四阶近似从下方逐步接近精确动能}{\begin{tikzpicture}
\begin{scope}[x=6cm,y=4cm]\draw[axis](0,0)--(.94,0)node[right]{$\beta=v/c$};\draw[axis](0,0)--(0,.77)node[above]{$K/(m_0c^2)$};\draw[world,domain=0:.8,samples=70]plot(\x,{1/sqrt(1-(\x)^2)-1});\draw[blue,thick,domain=0:.8,samples=70]plot(\x,{.5*(\x)^2+.375*(\x)^4});\draw[red,dashed,domain=0:.8,samples=70]plot(\x,{.5*(\x)^2});\draw[guide](.8,0)--(.8,.666667);\node[below]at(.8,0){$0.8$};\node[right]at(.8,.666667){精确值};\node[blue,right]at(.8,.4736){四阶};\node[red,right]at(.8,.32){二阶};\end{scope}
\end{tikzpicture}}



\par\medskip\Needspace{4\baselineskip}\noindent\textbf{以精确动能为基准计算相对误差}\par\nobreak
图\ref{c3e5s}将减去静能后的精确动能、二阶近似和四阶近似画在同一坐标系。小$\beta$时三线靠近，速度增大后经典二阶式逐渐低于精确曲线。记$K_{\rm cl}=m_0v^2/2$，对非零的小速度，先将分子、分母同除以$K_{\rm cl}$：
\[
\frac{K-K_{\rm cl}}{K}
=\frac{\frac34\beta^2+O(\beta^4)}
       {1+\frac34\beta^2+O(\beta^4)}
=\frac34\beta^2+O(\beta^4).
\]
因此经典式低估动能，且相对误差首项为$3v^2/(4c^2)$。这里始终使用不随速度变化的静质量$m_0$；若以经典动能作分母，首项相同，更高阶项则应重新计算。
\end{solution}

\begin{question}
\begin{example}[电偶极子的远场]
$-q$ 与 $+q$ 分别位于 $z=-d/2$ 与 $z=d/2$。观察点 $P$ 的球坐标为 $(r,\theta)$，$\theta$ 从 $+z$ 轴起量，$r\gg d$。求远场电势、电场，并给出轴向和赤道方向的结果。
\end{example}
\fig{c3e6q}{两电荷距离相差一个一阶小量}{\begin{tikzpicture}
\draw[axis](0,-1.6)--(0,2.3)node[above]{$z$};\fill(0,.6)circle(2pt)node[left]{$+q$};\fill(0,-.6)circle(2pt)node[left]{$-q$};
\coordinate(P)at(3,1.5);\fill(P)circle(2pt)node[right]{$P$};\draw[world](0,0)--(P)node[midway,below]{$r$};
\draw[guide](0,.6)--(P);\draw[guide](0,-.6)--(P);\draw(0,.5)arc(90:26.565:.5)node[midway,above right]{$\theta$};
\draw[dim](-.7,-.6)--(-.7,.6)node[midway,left]{$d$};
\end{tikzpicture}}
\end{question}
\begin{solution}\textbf{分别展开两个电荷到场点的距离}\par\nobreak
取 $\mathcal K=1/(4\pi\varepsilon_0)$，偶极矩大小 $p=qd$。图 \ref{c3e6q} 中两距离满足
$r_\pm^2=r^2+d^2/4\mp rd\cos\theta$。以 $d/r$ 展开，
\[
 \frac1{r_\pm}
 =\frac1r\left(1\pm\frac{d\cos\theta}{2r}
 +O(d^2/r^2)\right).
\]
这里先将距离平方中的 $r^2$ 提出，再对
$(1+\varepsilon_\pm)^{-1/2}$ 展开，其中
$\varepsilon_\pm=d^2/(4r^2)\mp(d/r)\cos\theta$。
一阶项是 $-\varepsilon_\pm/2$，其中关于 $d$ 的一次项分别为 $\pm d\cos\theta/(2r)$，所以得到上面的倒数。相减时 $1/r$ 抵消，两项一次项相加为 $d\cos\theta/r^2$。由交换两电荷等价于 $d\to-d$，电势差为奇函数，全部偶次项也抵消，下一项从三次开始，故
\[
 V=\mathcal Kq(1/r_+-1/r_-)
 =\frac{\mathcal Kp\cos\theta}{r^2}
 +O(\mathcal Kqd^3/r^4).
\]
\fig{c3e6s}{分别固定角度与半径，用线段和圆弧计算电势变化}{\begin{tikzpicture}
\draw[axis](0,-.4)--(0,3.5)node[above]{$z$};\coordinate(P)at(30:2.8);\coordinate(R)at(30:3.8);\coordinate(T)at(5:2.8);
\draw[guide](0,0)--(R);\draw[guide](0,0)--(T);\node[below left]at(0,0){$O$};
\draw[world](0,0)--(P)node[midway,above left]{$r$};
\draw[flow,blue](P)--(R)node[midway,above]{$\dd r$};\draw[flow,red](P)arc(30:5:2.8)node[midway,right,xshift=8pt]{$r\dd\theta$};
\draw(0,.7)arc(90:30:.7);\node at(60:1){$\theta$};\draw[flow](30:1.3)arc(30:5:1.3);\node at(17.5:1.8){$\dd\theta$};
\fill(P)circle(1.5pt)node[above left]{$P$};\fill(R)circle(1.5pt)node[right]{$P_r$};\fill(T)circle(1.5pt)node[below right]{$P_\theta$};
\node at(2,-.65){$P\to P_r:\ \theta$ 不变；$P\to P_\theta:\ r$ 不变};
\end{tikzpicture}}



\par\medskip\Needspace{4\baselineskip}\noindent\textbf{由两个方向的位移读取电势梯度}\par\nobreak
按图 \ref{c3e6s}，$P\to P_r$ 的径向小位移为 $\dd r\,\bm e_r$，$P\to P_\theta$ 的增角小位移为 $r\dd\theta\,\bm e_\theta$。对 $r$ 求偏导时固定 $\theta$，对 $\theta$ 求偏导时固定 $r$，分别对应图中的线段和圆弧。

电势微分满足
$\dd V=-E_r\dd r-E_\theta r\dd\theta$，
与 $\dd V=(\partial V/\partial r)\dd r+(\partial V/\partial\theta)\dd\theta$ 比较系数，得到
\begin{align*}
E_r&=-\frac{\partial V}{\partial r}
\simeq-\mathcal Kp\cos\theta\,\frac{\dd(r^{-2})}{\dd r}\\
&=-\mathcal Kp\cos\theta(-2r^{-3})
=\frac{2\mathcal Kp\cos\theta}{r^3},\\
E_\theta&=-\frac1r\frac{\partial V}{\partial\theta}
\simeq-\frac{\mathcal Kp}{r^3}\frac{\dd(\cos\theta)}{\dd\theta}\\
&=-\frac{\mathcal Kp}{r^3}(-\sin\theta)
=\frac{\mathcal Kp\sin\theta}{r^3}.
\end{align*}
横向的 $1/r$ 来自图中实际弧长 $r\dd\theta$，不能遗漏。

\par\medskip\Needspace{4\baselineskip}\noindent\textbf{检查轴向与赤道面的场方向}\par\nobreak
轴正向 $\theta=0$ 时 $E_\theta=0$，$E_r=2\mathcal Kp/r^3$，径向正是 $+z$；赤道面 $\theta=\pi/2$ 时 $E_r=0$，$E_\theta=\mathcal Kp/r^3$，增角方向则为 $-z$。对试探电荷 $Q$，受力为 $Q\bm E$；$Q<0$ 时方向反转。

\par\medskip\Needspace{4\baselineskip}\noindent\textbf{合成为矢量并投影到固定坐标轴}\par\nobreak
把已求的两个场分量合成为矢量，可以避免在轴线与赤道面之间更换符号。图\ref{c3e6resolve}取偶极轴为 $+z$，有
$\bm e_z=\cos\theta\bm e_r-\sin\theta\bm e_\theta$。记$\bm p=p\bm e_z$，则$\bm p\cdot\bm e_r=p\cos\theta$。把$-\bm p=-p\cos\theta\bm e_r+p\sin\theta\bm e_\theta$代入，可将两个分量合并成
\begin{align*}
\bm E&=\frac{\mathcal Kp}{r^3}
 (2\cos\theta\bm e_r+\sin\theta\bm e_\theta)\\
&=\frac{\mathcal K}{r^3}
 [3(\bm p\cdot\bm e_r)\bm e_r-\bm p].
\end{align*}
\fig{c3e6resolve}{径向与角向分量的正方向}{\begin{tikzpicture}
\draw[axis](0,-.5)--(0,3.3)node[above]{$z$};
\coordinate(P)at(2.2,1.65);\draw[guide](0,0)--(P);
\draw[flow,blue](P)--++(.96,.72)node[right]{$\bm e_r$};
\draw[flow,red](P)--++(.72,-.96)node[right]{$\bm e_\theta$};
\fill(P)circle(1.3pt)node[above left]{$P$};\node[below left]at(0,0){$O$};
\draw(0,.7)arc(90:36.87:.7)node[midway,above right]{$\theta$};
\end{tikzpicture}}
再取图示子午面内垂直偶极轴的方向为 $x$，则
$\bm e_r=(\sin\theta,\cos\theta)$、
$\bm e_\theta=(\cos\theta,-\sin\theta)$。将两个分量投影，得到
\begin{align*}
E_x&=E_r\sin\theta+E_\theta\cos\theta
\\&=\frac{\mathcal Kp}{r^3}(2\sin\theta\cos\theta+\sin\theta\cos\theta)
=\frac{3\mathcal Kp}{r^3}\sin\theta\cos\theta,\\
E_z&=E_r\cos\theta-E_\theta\sin\theta
\\&=\frac{\mathcal Kp}{r^3}(2\cos^2\theta-\sin^2\theta)
=\frac{\mathcal Kp}{r^3}(3\cos^2\theta-1).
\end{align*}
\par\medskip\Needspace{4\baselineskip}\noindent\textbf{用正交分量比较同一球面上的场强}\par\nobreak
在 $\cos^2\theta=1/3$处，$E_z$虽为零，$E_x$却不为零；某一分量的抵消不能误判为总场消失。利用径向与角向正交，还得
\[|\bm E|=\sqrt{E_r^2+E_\theta^2}
 =\frac{\mathcal Kp}{r^3}\sqrt{1+3\cos^2\theta}.\]
因此同一远场球面上的场强大小从赤道的 $\mathcal Kp/r^3$连续变到轴线的 $2\mathcal Kp/r^3$。这些结果均限于 $r$远大于电荷间距，不能把偶极近似延伸到两个电荷附近。
\end{solution}

\begin{question}
\begin{example}[指数密度液体中的浮杆]
液面以下深度为 $z$ 处的密度为 $\rho(z)=\rho_0\e^{z/\ell}$。竖直刚性杆长 $\ell$、截面积 $S$、密度 $3\rho_0/2$。液体静止，杆保持竖直。求平衡时杆顶相对液面的高度。
\end{example}
\fig{c3e7q}{深度坐标从自由液面向下}{\begin{tikzpicture}
\fill[blue!7](-2,-2.5)rectangle(2,0);\draw[world](-2,0)--(2,0)node[right]{液面};
\draw[axis](-1.5,.5)--(-1.5,-2.8)node[below]{$z$};
\draw[fill=gray!15](-.3,.3)rectangle(.3,-2.1);\draw[dim](.8,.3)--(.8,-2.1)node[midway,right]{$\ell$};
\node[left]at(-1.5,0){$0$};
\end{tikzpicture}}
\end{question}
\begin{solution}\textbf{先假设部分浸没，按深度积分压强}\par\nobreak
先假定图 \ref{c3e7q} 所示的部分浸没状态，令杆底深度为 $b$。静水压梯度为 $\dd p/\dd z=\rho(z)g$，顶面在空气中，故净浮力等于
\[
 \begin{aligned}
 F_B&=S[p(b)-p(0)]
=\rho_0gS\int_0^b\e^{z/\ell}\dd z\\
&=\rho_0gS[\ell\e^{z/\ell}]_0^b
=\rho_0gS(\ell\e^{b/\ell}-\ell\e^0)\\
&=\rho_0gS\ell(\e^{b/\ell}-1).
 \end{aligned}
\]
指数项的原函数可由$\dd(\ell\e^{z/\ell})/\dd z=\ell\e^{z/\ell}(1/\ell)=\e^{z/\ell}$核验。令浮力等于杆重$3\rho_0gS\ell/2$，约去正的共同因子$\rho_0gS\ell$，得到$\e^{b/\ell}-1=3/2$，所以$\e^{b/\ell}=5/2$，取对数得
\[
 b=\ell\ln(5/2)\simeq0.91629\ell<\ell.
\]
\fig{c3e7s}{平衡后顶端露出液面，底端位于深度 $b$}{\begin{tikzpicture}
\draw[world](-1.7,0)--(1.8,0);\draw[fill=gray!15](-.3,.22)rectangle(.3,-2.42);
\draw[dim](.8,0)--(.8,-2.42)node[midway,right]{$b=\ell\ln(5/2)$};
\draw[force,blue](0,-1.1)--(0,-.2)node[left]{$F_B$};\draw[force,red](0,-1.1)--(0,-2)node[left]{$mg$};
\node[above]at(0,.35){露出 $\ell-b$};
\end{tikzpicture}}



\par\medskip\Needspace{4\baselineskip}\noindent\textbf{核验浸没状态并排除不成立的支路}\par\nobreak
图 \ref{c3e7s} 中假设得到验证，顶端露出高度为
\[
 \boxed{h_{\rm top}=\ell[1-\ln(5/2)]\simeq0.08371\ell}.
\]
若强行假定完全浸没，杆顶深度为$z_0\ge0$，杆底深度为$z_0+\ell$，上下边界必须同时改变：
\begin{align*}
F_B&=\rho_0gS\int_{z_0}^{z_0+\ell}\e^{z/\ell}\dd z
=\rho_0gS[\ell\e^{z/\ell}]_{z_0}^{z_0+\ell}\\
&=\rho_0gS\ell
 [\e^{(z_0+\ell)/\ell}-\e^{z_0/\ell}]
=\rho_0gS\ell\e^{z_0/\ell}(\e-1).
\end{align*}

与重量相等、提出 $\e^{z_0/\ell}$ 后，应满足
$\e^{z_0/\ell}=\tfrac{3}{2(\e-1)}<1$，反而要求 $z_0<0$，不自洽。对实际部分浸没状态，
\[
\frac{\dd F_B}{\dd b}
=\rho_0gS\ell\e^{b/\ell}\frac1\ell
=\rho_0gS\e^{b/\ell}>0.
\]
下沉小量 $\delta b>0$ 时重量不变、浮力增加，向下的合力变化为 $-\rho_0gS\e^{b/\ell}\delta b<0$，故竖直平移方向稳定。


\par\medskip\Needspace{4\baselineskip}\noindent\textbf{在部分浸没区间构造有效势能}\par\nobreak
% thirty-c3-floatstable
还可将上述回复力写成势能，检查局部判稳与有限位移是否一致。令平衡底深为 $b_0=\ell\ln(5/2)$，杆整体向下位移 $\eta$，则底深为 $b_0+\eta$。在仍部分浸没的区间 $-b_0<\eta<\ell-b_0$，净向下力为
\begin{align*}
Q(\eta)&=mg-F_B(b_0+\eta)\\
 &=\rho_0gS\ell\left[\frac32-\left(\frac52\e^{\eta/\ell}-1\right)\right]
 =\frac52\rho_0gS\ell(1-\e^{\eta/\ell}).
\end{align*}
将浮力和重力合并为位置势能，$\dd U/\dd\eta=-Q$。用 $\xi$ 作积分变量，从平衡位置到 $\eta$ 积分，先写
\[
 U(\eta)-U(0)=\frac52\rho_0gS\ell
 \int_0^\eta(\e^{\xi/\ell}-1)\dd\xi
 =\frac52\rho_0gS\ell[\ell\e^{\xi/\ell}-\xi]_0^\eta.
\]
指数项可令 $u=\xi/\ell$、$\dd\xi=\ell\dd u$，其原函数才是 $\ell\e^{\xi/\ell}$。代入上下限，得到
\[U(\eta)-U(0)=\frac52\rho_0gS\ell
 \left[\ell(\e^{\eta/\ell}-1)-\eta\right].\]
指数函数满足 $\e^\xi\ge1+\xi$，因此方括号非负，且仅在 $\eta=0$为零。

图\ref{c3e7potential}显示平衡位置确为势能低点，而不仅是合力为零的点。
\fig{c3e7potential}{部分浸没范围内的有效势能具有严格极小值}{\begin{tikzpicture}
\draw[axis](-2,0)--(2.2,0)node[right]{$\xi=\eta/\ell$};
\draw[axis](0,-.25)--(0,2.3)node[above]{$\Delta U/U_*$};
\draw[world,blue]plot[domain=-.08:.08,samples=70]({20*\x},{500*(\x*\x/2+\x*\x*\x/6+\x*\x*\x*\x/24+\x*\x*\x*\x*\x/120+\x*\x*\x*\x*\x*\x/720)});
\node[below]at(-1.6,0){$-.08$};\node[below]at(1.6,0){$.08$};
\node[above left]at(0,0){$0$};
\node at(0,-.9){按六阶展开放大显示，$U_*=\frac52\rho_0gS\ell^2$};
\end{tikzpicture}}

\par\medskip\Needspace{4\baselineskip}\noindent\textbf{核对势能展开与回复刚度}\par\nobreak
继续展开 $\e^{\eta/\ell}=1+\eta/\ell+\eta^2/(2\ell^2)+O(\eta^3/\ell^3)$，线性项抵消，故
\[U(\eta)-U(0)=\frac54\rho_0gS\eta^2
 +O\!\left(\rho_0gS\frac{|\eta|^3}{\ell}\right).\]
二次项系数为正，给出回复刚度 $5\rho_0gS/2$，与前面的浮力导数相同。这里检验竖直平移的静态稳定性；计算实际振动频率时，还需明确液体运动和阻尼的处理条件。

\end{solution}

\begin{question}
\begin{example}[抛物线与椭圆的曲率]
求抛物线 $y=ax^2$ 在顶点的曲率半径，以及椭圆 $x^2/A^2+y^2/B^2=1$ 在两类轴端点的曲率半径，$a,A,B>0$。
\end{example}
\fig{c3e8q}{椭圆两类端点的弯曲程度不同}{\begin{tikzpicture}
\draw[world](0,0)ellipse(2.2 and 1.2);\draw[axis](-2.6,0)--(2.7,0)node[right]{$x$};\draw[axis](0,-1.5)--(0,1.8)node[above]{$y$};
\fill(2.2,0)circle(2pt);\node[right]at(2.6,-.65){$P(A,0)$};\fill(0,1.2)circle(2pt)node[above right]{$Q(0,B)$};
\end{tikzpicture}}
\end{question}
\begin{solution}\textbf{分别选取函数式与椭圆参数式}\par\nobreak
对抛物线固定$a$，先在一般位置求导：$y'=a\,2x=2ax$，$y''=2a$。顶点$x=0$处才有$y'=0$，代入曲率公式得到$\rho=(1+0^2)^{3/2}/|2a|=1/(2a)$。其中$a$的量纲为逆长度，所得半径量纲正确。

对图 \ref{c3e8q} 的椭圆，取正则参数 $x=A\cos\theta,y=B\sin\theta$。有
\[
 x'=-A\sin\theta,\quad y'=B\cos\theta,\quad
 x''=-A\cos\theta,\quad y''=-B\sin\theta,
\]
撇号在这里表示对参数 $\theta$ 求导。代入行列式，
$x'y''-y'x''=AB\sin^2\theta+AB\cos^2\theta=AB$。
同时 $x'^2+y'^2=A^2\sin^2\theta+B^2\cos^2\theta$，两部分都代入曲率公式。
\fig{c3e8s}{密切圆半径沿内法线}{\begin{tikzpicture}
\draw[world](0,0)ellipse(2.2 and 1.2);\draw[force,blue](2.2,0)--(1.5455,0)node[midway,above]{$\rho_P$};
\draw[force,red](0,1.2)--(0,-2.8333)node[midway,left]{$\rho_Q$};
\draw[guide](2.2,-.7)--(2.2,.7);\draw[guide](-.8,1.2)--(.8,1.2);
\end{tikzpicture}}

\[
 \rho(\theta)=
 \frac{(A^2\sin^2\theta+B^2\cos^2\theta)^{3/2}}{AB}.
\]


\par\medskip\Needspace{4\baselineskip}\noindent\textbf{在轴端点代入并比较弯曲程度}\par\nobreak
在 $P$ 处 $\theta=0$，分子为 $(B^2)^{3/2}=B^3$，除以 $AB$ 得 $B^2/A$；在 $Q$ 处 $\theta=\pi/2$，分子变为 $A^3$，除以 $AB$ 得 $A^2/B$。因此
$\boxed{\rho_P=B^2/A,\quad \rho_Q=A^2/B}$。
图 \ref{c3e8s} 的箭头表示内法线方向；当 $A>B$ 时长轴端更尖，半径更小。$A=B=R$ 时两式同时恢复圆的半径 $R$。


\par\medskip\Needspace{4\baselineskip}\noindent\textbf{沿内法线确定一般位置的曲率中心}\par\nobreak
% thirty-c3-ellipse
已知半径后，还应能在图上确定曲率中心。记
$D=A^2\sin^2\theta+B^2\cos^2\theta$。沿参数增加方向的单位切向为
$\bm e_t=(-A\sin\theta,B\cos\theta)/\sqrt D$；将其逆时针转 $90^\circ$，得到指向椭圆内部的单位法向
\[\bm e_n=\frac{(-B\cos\theta,-A\sin\theta)}{\sqrt D}.\]
所以曲率中心 $C$的坐标应由 $\bm r_C=\bm r_P+\rho\bm e_n$确定，而不是任意沿半径 $OP$取点。代入已求的 $\rho=D^{3/2}/(AB)$，分别整理两个分量：
\begin{align*}
x_C&=A\cos\theta-\frac D A\cos\theta
 =\frac{A^2-B^2}{A}\cos^3\theta,\\
y_C&=B\sin\theta-\frac D B\sin\theta
 =\frac{B^2-A^2}{B}\sin^3\theta.
\end{align*}
\fig{c3e8center}{一般位置的曲率中心位于内法线上}{\begin{tikzpicture}
\draw[world](0,0)ellipse(2.2 and 1.2);
\draw[axis](-2.6,0)--(2.7,0)node[right]{$x$};
\draw[axis](0,-1.65)--(0,1.8)node[above]{$y$};
\coordinate(P)at(1.55563,.84853);\coordinate(C)at(.54640,-1.00173);
\path (P);\draw[flow,blue](P)--(C)node[midway,right,xshift=10pt,yshift=-8pt]{$\rho\bm e_n$};
\draw[flow,red](P)--++(-.9,.49)node[above]{$\bm e_t$};
\fill(P)circle(1.5pt)node[above right]{$P$};\fill(C)circle(1.5pt)node[below]{$C$};
\node[below left]at(0,0){$O$};
\end{tikzpicture}}
图\ref{c3e8center}中的 $PC$与切向垂直，长度为 $\rho$；一般位置的法线并不经过 $O$。在长轴端 $P=(A,0)$，曲率中心为 $((A^2-B^2)/A,0)$，两点之差的长度正是 $B^2/A$。在短轴端，曲率中心可能落在椭圆之外；这不影响它作为局部密切圆圆心的意义。

若令 $A=B=R$，则所有曲率中心均化为 $O$、所有半径均化为 $R$，同时检验了位置和大小。求几何曲率时可以自由选取正则参数；真实速度大小仍由时间规律决定，不能由这个参数化直接认定物体做匀速运动。
\end{solution}

\begin{question}
\begin{example}[移动切线的切点]
固定圆半径 $R$、圆心 $O$。点 $A=(x,0)$ 位于圆外，$x>R$，以速度 $\dot x=u>0$ 向右运动。从 $A$ 作上侧切线，切点为 $P$。求 $P$ 的坐标、速度大小以及切线段 $AP$ 长度的变化率。
\end{example}
\fig{c3e9q}{半径与切线垂直}{\begin{tikzpicture}[scale=1.15]
\draw[world](0,0)circle(1.2);\coordinate(A)at(3.2,0);\coordinate(P)at(.45,1.1124);
\draw[axis](-1.5,0)--(4,0)node[right]{$x$};\draw[world](0,0)--(P)--(A);
\fill(P)circle(2pt)node[above]{$P$};\fill(A)circle(2pt)node[below]{$A$};\node[below]at(0,0){$O$};
\draw[flow](3.2,.65)--(4,.65)node[above]{$u$};
\end{tikzpicture}}
\end{question}
\begin{solution}\textbf{联立切线垂直条件和圆方程}\par\nobreak
图 \ref{c3e9q} 中设 $P=(X,Y)$。半径与切线垂直，所以
$\overrightarrow{OP}\cdot\overrightarrow{AP}=X(X-x)+Y^2=0$。
又因 $X^2+Y^2=R^2$，有 $R^2-xX=0$，从而 $X=R^2/x$。代回圆方程，并由上侧切点的 $Y>0$ 取正根，
\[
 X=\frac{R^2}{x},\qquad Y=\frac{R\sqrt{x^2-R^2}}x.
\]
求导时$\dot x=u$，而$X,Y$都通过$x(t)$随时间变化。先记$W=\sqrt{x^2-R^2}>0$，按链式法则求出根式的导数，再对$Y=RW/x$使用商法则：
\begin{align*}
W'&=\frac12(x^2-R^2)^{-1/2}(2x)=\frac xW,\\
\frac{\dd Y}{\dd x}
&=R\frac{W'x-W}{x^2}
=\frac{R(x^2-W^2)}{x^2W}
=\frac{R[x^2-(x^2-R^2)]}{x^2W}\\
&=\frac{R^3}{x^2\sqrt{x^2-R^2}}.
\end{align*}
再乘以 $\dot x=u$，得到
\[
 \dot X=-\frac{R^2u}{x^2},\qquad
 \dot Y=\frac{R^3u}{x^2\sqrt{x^2-R^2}}.
\]
\fig{c3e9s}{切点沿圆运动，而非沿切线滑向外点}{\begin{tikzpicture}[scale=1.15]
\draw[world](0,0)circle(1.2);\coordinate(P)at(.45,1.1124);\draw[guide](0,0)--(P);
\draw[force,blue](P)--++(-.9,.364)node[above]{$\bm v_P$};\draw[guide](P)--(3.2,0);\node[right]at(3.2,0){$A$};\node[below]at(0,0){$O$};
\end{tikzpicture}}



\par\medskip\Needspace{4\baselineskip}\noindent\textbf{合成分量速度并检查运动方向}\par\nobreak
图 \ref{c3e9s} 的两个分量一负一正，切点向左上方移动。计算速率平方时先提出公因子：
\[
 v_P^2=\frac{R^4u^2}{x^4}
 \left(1+\frac{R^2}{x^2-R^2}\right)
 =\frac{R^4u^2}{x^2(x^2-R^2)}.
\]
用 $x>R,u>0$ 取正根，
\[
 \boxed{v_P=\frac{R^2u}{x\sqrt{x^2-R^2}}}.
\]
由直角三角形$AP=\sqrt{x^2-R^2}$，再次追踪$x(t)$的内层导数：
\[
\frac{\dd(AP)}{\dd t}
=\frac{1}{2\sqrt{x^2-R^2}}\,2x\dot x
=\frac{xu}{\sqrt{x^2-R^2}}.
\]
这是切线段的伸长率，与切点速率不同。$x\to R^+$ 时二者发散，提醒我们外点接近圆周时，切点几何对应关系对位置变化异常敏感。
\end{solution}

\section{本讲小结}
极值分析要从定义域、候选点和导数符号展开，势能判稳则必须明确允许的扰动方向。近似展开先找无量纲小参数，再保留足以决定结果的最低非零阶；发生相减时应特别检查低阶项是否抵消。曲率半径描述路径的局部几何，几何点的速度则还取决于驱动条件，二者需要不同的信息。
\section{练习}
\begin{question}
\begin{drill}[带误差的近似]\label{c3d1}
用二阶展开估算 $\sqrt{4.01}$，并给出误差上界。再估计把 $\sin\theta$ 近似为 $\theta$ 时，在 $|\theta|\le0.1$ 内相对 $\theta$ 的误差上界。
\end{drill}
\fig{c3d1q}{展开点与实际求值点相邻}{\begin{tikzpicture}
\draw[axis](0,0)--(5.5,0)node[right]{$x$};\node[below]at(1,0){$4$};\node[below]at(4,0){$4.01$};\draw[dim](1,.5)--(4,.5)node[midway,above]{$h=0.01$};
\end{tikzpicture}}\answerlink{c3d1}
\end{question}
\begin{question}
\begin{drill}[分层大气中的圆形光路]\label{c3d2}
球对称大气折射率为 $n(r)=1+\delta\e^{-(r-r_0)/b}$，$r_0=6400\,\mathrm{km}$，$b=8772\,\mathrm m$。已知近地层实际折射率为 $1.0003$。若光线要沿 $r=r_0$ 的圆周传播，所需密度约为实际密度的多少倍？假定 $n-1$ 与密度成正比，并使用分层介质光线曲率关系 $\kappa=-n'(r)/n(r)$。
\end{drill}
\fig{c3d2q}{折射率梯度朝向地心}{\begin{tikzpicture}[scale=.9]
\draw[world](0,0)circle(1.4);\draw[guide](0,0)circle(1.8);\draw[ray,blue](1.8,0)arc(0:95:1.8);
\draw[flow](40:2.2)--(40:1.5)node[right]{$\nabla n$};\node[below]at(0,0){$O$};\draw[guide](0,0)--(0,1.8)node[midway,right]{$r_0$};
\end{tikzpicture}}\answerlink{c3d2}
\end{question}
\begin{question}
\begin{drill}[闭区间的最大与最小]\label{c3d3}
求 $f(x)=x^3-3x$ 在 $[-2,3]$ 上的最大值和最小值，说明为何只解 $f'=0$ 不足以完成此题。
\end{drill}
\fig{c3d3q}{端点也属于可取位置}{\begin{tikzpicture}
\draw[axis](-2.7,0)--(3.7,0)node[right]{$x$};\draw[very thick,blue](-2,0)--(3,0);\foreach \x in {-2,-1,1,3}{\fill(\x,0)circle(2pt)node[below]{$\x$};}
\end{tikzpicture}}\answerlink{c3d3}
\end{question}
\begin{question}
\begin{drill}[曲率与参数选择]\label{c3d4}
质点轨迹为 $x=A\cos\omega t,\ y=B\sin\omega t$，$A>B>0,\omega>0$。求 $t=0$ 时速率、加速度和曲率半径；再求 $\omega t=\pi/4$ 时的切向加速度，判断它是否一般为零。
\end{drill}
\fig{c3d4q}{恒定参数角速度不等于恒定速率}{\begin{tikzpicture}
\draw[world](0,0)ellipse(2 and 1);\fill(2,0)circle(2pt)node[right]{$t=0$};\fill(1.4142,.7071)circle(2pt)node[above right]{$\omega t=\pi/4$};
\draw[guide](0,0)--(1.4142,.7071);\node[below]at(0,0){$O$};
\end{tikzpicture}}\answerlink{c3d4}
\end{question}

\LectureAnswerSection

\begin{drillsolution}{c3d1}\textbf{在已知点附近展开根式并估计余项}\par\nobreak
图\ref{c3d1q}中$h=0.01$。把$f(x)=\sqrt x$写成$x^{1/2}$，逐次应用幂函数求导法则：
\begin{align*}
f'(x)&=\frac12x^{-1/2}=\frac1{2\sqrt x},\\
f''(x)&=\frac12\left(-\frac12\right)x^{-3/2}
=-\frac1{4x^{3/2}},\\
f'''(x)&=-\frac14\left(-\frac32\right)x^{-5/2}
=\frac3{8x^{5/2}}.
\end{align*}
在四处代入得到 $f(4)=2$、$f'(4)=1/4$、$f''(4)=-1/32$，所以二阶展开为
\[
 \sqrt{4+h}=2+\frac h4-\frac{h^2}{64}+R_2,
\]
代入$h=0.01$，依次计算$h/4=0.0025$、$h^2/64=0.0000015625$，得到$\sqrt{4.01}\simeq2.0024984375$。前面已求得$f'''(x)=3/(8x^{5/2})$；它在正半轴随$x$增加而减小，所以在$[4,4.01]$上的最大值为$f'''(4)=3/(8\cdot32)=3/256$。按拉格朗日余项，
$|R_2|\le(3/256)h^3/3!=h^3/512=1.953125\times10^{-9}$。
\fig{c3d1s}{放大余项与其严格上界之比，整个区间都在一以下}{\begin{tikzpicture}
\draw[axis](0,0)--(4.8,0)node[right]{$h$};\draw[axis](0,0)--(0,2.6)node[above]{$R_2/(h^3/512)$};\fill[blue!7](0,0)rectangle(4,2);\draw[red,dashed](0,2)--(4.2,2)node[right]{上界 $1$};\draw[world]plot coordinates{(0.00000,2.0000000) (0.05000,1.9804692) (0.10000,1.9609392) (0.15000,1.9414101) (0.20000,1.9218818) (0.25000,1.9023544) (0.30000,1.8828279) (0.35000,1.8633022) (0.40000,1.8437773) (0.45000,1.8242533) (0.50000,1.8047302) (0.55000,1.7852079) (0.60000,1.7656865) (0.65000,1.7461659) (0.70000,1.7266462) (0.75000,1.7071273) (0.80000,1.6876093) (0.85000,1.6680922) (0.90000,1.6485759) (0.95000,1.6290604) (1.00000,1.6095458) (1.05000,1.5900321) (1.10000,1.5705192) (1.15000,1.5510071) (1.20000,1.5314960) (1.25000,1.5119856) (1.30000,1.4924761) (1.35000,1.4729675) (1.40000,1.4534597) (1.45000,1.4339528) (1.50000,1.4144468) (1.55000,1.3949415) (1.60000,1.3754372) (1.65000,1.3559337) (1.70000,1.3364310) (1.75000,1.3169292) (1.80000,1.2974282) (1.85000,1.2779281) (1.90000,1.2584289) (1.95000,1.2389305) (2.00000,1.2194330) (2.05000,1.1999363) (2.10000,1.1804404) (2.15000,1.1609454) (2.20000,1.1414513) (2.25000,1.1219580) (2.30000,1.1024656) (2.35000,1.0829740) (2.40000,1.0634833) (2.45000,1.0439934) (2.50000,1.0245044) (2.55000,1.0050162) (2.60000,0.9855289) (2.65000,0.9660424) (2.70000,0.9465568) (2.75000,0.9270720) (2.80000,0.9075881) (2.85000,0.8881050) (2.90000,0.8686228) (2.95000,0.8491414) (3.00000,0.8296609) (3.05000,0.8101813) (3.10000,0.7907025) (3.15000,0.7712245) (3.20000,0.7517474) (3.25000,0.7322711) (3.30000,0.7127957) (3.35000,0.6933212) (3.40000,0.6738474) (3.45000,0.6543746) (3.50000,0.6349026) (3.55000,0.6154314) (3.60000,0.5959611) (3.65000,0.5764917) (3.70000,0.5570230) (3.75000,0.5375553) (3.80000,0.5180884) (3.85000,0.4986223) (3.90000,0.4791571) (3.95000,0.4596928) (4.00000,0.4402293)};
\node[left]at(0,0){$.998$};\node[left]at(0,1){$.999$};\node[left]at(0,2){$1.000$};\node[below]at(0,0){$0$};\node[below]at(4,0){$.01$};\node at(2,-.8){纵轴只显示 $[.998,1]$，$h=0$ 取极限};
\end{tikzpicture}}



\par\medskip\Needspace{4\baselineskip}\noindent\textbf{以严格余项界确定小角允许范围}\par\nobreak
图 \ref{c3d1s}放大显示余项与上界的比值：它为正且小于 $1$，所以实际值在二阶近似上方，但增量不超过 $h^3/512$。这个界来自三阶导数，不能用保留的小数位数替代。对正弦在零点保留到二阶，常数、二次项均为零，三阶导数为 $-\cos\theta$、绝对值不超过一，因此余项界为 $|\sin\theta-\theta|\le|\theta|^3/3!$。再除以 $|\theta|>0$，得到
$|\sin\theta-\theta|/|\theta|\le\theta^2/6\le1/600$；零点以极限补充为零。这里误差的比较基准是 $\theta$，不是 $\sin\theta$。
\end{drillsolution}
\begin{drillsolution}{c3d2}\textbf{沿内法线把圆轨道曲率代入折射率梯度}\par\nobreak
令 $\delta=n(r_0)-1>0$。先对题给指数逐层求导：$n'(r)=\delta\e^{-(r-r_0)/b}\,[-1/b]$。在图\ref{c3d2q}的地表位置代入 $r=r_0$，得到 $n(r_0)=1+\delta$、$n'(r_0)=-\delta/b$；光线向折射率较大的内层弯曲。半径 $r_0$的圆形光路要求内法线曲率为 $1/r_0$，代入题给曲率关系，
\[\frac1{r_0}=-\frac{n'(r_0)}{n(r_0)}
 =\frac{\delta}{b(1+\delta)}.\]
\fig{c3d2s}{局部切向传播时的内法线曲率}{\begin{tikzpicture}
\draw[world](0,0)arc(0:70:1.7);\coordinate(P)at(0,0);\draw[ray,blue](P)--(0,1.3)node[above]{切向};
\draw[force,red](P)--(-1,0)node[left]{内法线};\node[below right]at(P){$P$};
\end{tikzpicture}}



\par\medskip\Needspace{4\baselineskip}\noindent\textbf{解出折射率增量并换算密度}\par\nobreak
图\ref{c3d2s}标明内法线方向，所以曲率取正。将等式两边乘正数 $br_0(1+\delta)$，再移项：
\[b(1+\delta)=r_0\delta,\qquad
b=(r_0-b)\delta,\qquad
\delta=\frac b{r_0-b}.\]
$n-1$与密度成正比，比较的应是折射率增量，而非两个接近一的折射率本身。先把 $r_0=6400\,\mathrm{km}$换成 $6.4\times10^6\,\mathrm m$，再与实际增量0.0003相比，得到
\[\frac{\rho_{\rm req}}{\rho_{\rm act}}
 =\frac{\delta}{0.0003}
 =\frac{8772}{(6400000-8772)\,0.0003}
 \simeq\boxed{4.575}.\]
若在曲率式中预先将 $1+\delta$近似为一，会得到约4.57；完整表达式则明确保留了这一步近似的位置。


\par\medskip\Needspace{4\baselineskip}\noindent\textbf{用任意径向扰动独立核验光路条件}\par\nobreak
还可用光程驻值独立核验圆形光路的条件，其作用是解释曲率条件的来由。取极角 $\phi$为参数，将邻近路径写成 $r(\phi)=r_0+\epsilon\eta(\phi)$，其中 $\eta$为光滑的周期函数。路径微元为
\[\dd s=\sqrt{r^2+(\dd r/\dd\phi)^2}\dd\phi
 =[r_0+\epsilon\eta(\phi)]\dd\phi+O(\epsilon^2).\]
导数平方只贡献二阶小量；折射率的一阶展开为
\[n(r)=n(r_0)+\epsilon n'(r_0)\eta(\phi)+O(\epsilon^2).\]
两式相乘，沿一周积分，得到光程的一阶变化
\[\delta\mathcal L=\epsilon\int_0^{2\pi}
 [n(r_0)+r_0n'(r_0)]\eta(\phi)\dd\phi.\]
\fig{c3d2variation}{路径略向外移时，长度增加而折射率减小}{\begin{tikzpicture}
\draw[blue,world](0,0)circle(1.5);\draw[red,dashed](0,0)circle(1.9);
\draw[guide](0,0)--(35:1.9);\draw[dim](35:1.5)--(35:1.9)node[right]{$\epsilon\eta$};
\node at(0,-2.4){蓝：$r_0$\quad 红：邻近路径};
\node[below left]at(0,0){$O$};
\end{tikzpicture}}
图\ref{c3d2variation}中路径向外移动时，几何长度增加、折射率减小。

要使任意小扰动下的一阶变化为零，二者须满足 $n(r_0)+r_0n'(r_0)=0$，即 $-n'(r_0)/n(r_0)=1/r_0$，与前面使用的曲率条件相同。两种方法由此共同核验了圆形光路的条件和弯曲方向；候选光路的扰动稳定性还需进一步检查二阶变化。
\end{drillsolution}
\begin{drillsolution}{c3d3}\textbf{列出驻点、端点并逐一比较函数值}\par\nobreak
图\ref{c3d3q}给出的 $f(x)=x^3-3x$在闭区间 $[-2,3]$连续，故全局最值一定存在。先求导并解出内部驻点：
\[f'(x)=3x^2-3=3(x-1)(x+1)=0,
\qquad x=-1,1.\]


\par\medskip\Needspace{4\baselineskip}\noindent\textbf{按单调区间收集全部候选点并比较}\par\nobreak
导数在 $[-2,-1]$为正，在 $[-1,1]$为负，在 $[1,3]$为正；这三个单调区间覆盖题设区间，所以只需比较它们的边界 $-2,-1,1,3$。逐项代入原函数，
\[f(-2)=-8+6=-2,\quad f(-1)=-1+3=2,\]
\[f(1)=1-3=-2,\qquad f(3)=27-9=18.\]
\fig{c3d3s}{端点与驻点都要比较，局部极大不保证全局最大}{\begin{tikzpicture}[x=1cm,y=.13cm]
\draw[->](-2.7,0)--(3.8,0)node[right]{$x$};\draw[->](-2.7,-9)--(-2.7,21)node[above]{$f(x)$};
\foreach \x/\y in {-2/-2,-1/2,1/-2,3/18}{\draw[blue,line width=5pt](\x,0)--(\x,\y);\fill(\x,\y)circle(1.5pt);\node[below]at(\x,-7){$\x$};}
\node[below right]at(-2,-2){$-2$};\node[below right]at(1,-2){$-2$};\node[above]at(-1,2){$2$};\node[above]at(3,18){$18$};
\end{tikzpicture}}
由图\ref{c3d3s}的四个候选值，最大值为 $18$，在 $x=3$取得；最小值为 $-2$，在 $x=-2,1$两处取得。内部驻点 $x=-1$虽由增转减，是局部极大点，其函数值2却小于右端点的18。端点无需满足 $f'=0$，因此只列驻点不能完成闭区间最值问题。
\end{drillsolution}
\begin{drillsolution}{c3d4}\textbf{利用轴端点的正交关系求曲率}\par\nobreak
先在图\ref{c3d4q}的固定直角坐标中求一般时刻的分量。$A,B,\omega$保持不变，每次时间求导都带上内层导数$\dd(\omega t)/\dd t=\omega$：
\begin{align*}
v_x&=A[-\sin(\omega t)]\omega=-A\omega\sin\omega t,\\
v_y&=B\cos(\omega t)\omega=B\omega\cos\omega t,\\
a_x&=-A\omega\cos(\omega t)\omega=-A\omega^2\cos\omega t,\\
a_y&=B\omega[-\sin(\omega t)]\omega=-B\omega^2\sin\omega t.
\end{align*}
现在代入$t=0$，有$\bm v=(0,B\omega)$、$v=B\omega$，$\bm a=(-A\omega^2,0)$、$|\bm a|=A\omega^2$。两矢量垂直，加速度全为法向，因此
\[
\rho=\frac{v^2}{|\bm a|}
=\frac{B^2\omega^2}{A\omega^2}=\frac{B^2}{A}.
\]
\fig{c3d4s}{轴端点可以直接用全部加速度求曲率}{\begin{tikzpicture}[scale=.85]
\fill(0,0)circle(2pt)node[below right]{$P(A,0)$};\draw[force,blue](0,0)--(0,1.5)node[above]{$B\omega$};\draw[force,red](0,0)--(-1.7,0)node[left]{$A\omega^2$};
\end{tikzpicture}}



\par\medskip\Needspace{4\baselineskip}\noindent\textbf{在一般位置分离切向与法向加速度}\par\nobreak
图\ref{c3d4s}的垂直关系只在该类特殊位置成立。对一般位置，使用上面逐分量求得的速度和加速度计算点积。
点积为
$\bm v\cdot\bm a=\omega^3(A^2-B^2)\sin\omega t\cos\omega t$，
速率为$\omega\sqrt{A^2\sin^2\omega t+B^2\cos^2\omega t}$。由于$v^2=\bm v\cdot\bm v$，两边对时间求导得到$2v\dot v=2\bm v\cdot\bm a$。在本题$v>0$，因此切向加速度$a_\tau=\dot v=(\bm v\cdot\bm a)/v$，即
\[
 a_\tau=\frac{\bm v\cdot\bm a}{v}
 =\frac{\omega^2(A^2-B^2)\sin\omega t\cos\omega t}
 {\sqrt{A^2\sin^2\omega t+B^2\cos^2\omega t}}.
\]
在 $\omega t=\pi/4$ 时，分子中的 $\sin\omega t\cos\omega t=1/2$，分母根式为 $\sqrt{(A^2+B^2)/2}$，化简后为
$\omega^2(A^2-B^2)/\sqrt{2(A^2+B^2)}>0$。
所以恒定参数角速度下，椭圆上的速率正在增加；只有 $A=B$ 的圆才处处无切向加速度。

% thirty-c3-ellipsecomponents

\par\medskip\Needspace{4\baselineskip}\noindent\textbf{由速度与加速度的叉积求法向分量}\par\nobreak
为补全一般位置的分解，令 $\phi=\omega t$、
$D=A^2\sin^2\phi+B^2\cos^2\phi$。平面运动的叉积仅有垂直纸面的分量，逐项计算为
\begin{align*}
(\bm v\times\bm a)_z
 &=(-A\omega\sin\phi)(-B\omega^2\sin\phi)
 -(B\omega\cos\phi)(-A\omega^2\cos\phi)\\
 &=AB\omega^3(\sin^2\phi+\cos^2\phi)=AB\omega^3.
\end{align*}
因切向部分与 $\bm v$平行，只有法向部分进入叉积。由 $|\bm v\times\bm a|=v a_n$，可得
\[a_n=\frac{AB\omega^2}{\sqrt D},\qquad
\rho=\frac{v^2}{a_n}=\frac{D^{3/2}}{AB}.\]
\fig{c3d4components}{一般位置先分解加速度，再用法向分量求曲率}{\begin{tikzpicture}
\coordinate(P)at(0,0);\coordinate(T)at(-.842,.459);
\coordinate(N)at(-.714,-1.308);\coordinate(A)at(-1.556,-.849);
\draw[guide](T)--(A)--(N);
\draw[flow,blue](P)--(T)node[above left]{$a_\tau\bm e_\tau$};
\draw[flow,red](P)--(N)node[below]{$a_n\bm e_n$};
\draw[flow](P)--(A)node[left]{$\bm a$};
\fill(P)circle(1.4pt)node[right]{$P$};
\end{tikzpicture}}
图\ref{c3d4components}的两分量相互垂直；总加速度的大小应满足
$|\bm a|^2=a_\tau^2+a_n^2$。在 $\phi=\pi/4$处，
\[a_n=\frac{\sqrt2AB\omega^2}{\sqrt{A^2+B^2}},\qquad
a_\tau=\frac{(A^2-B^2)\omega^2}{\sqrt{2(A^2+B^2)}}.\]
两式平方相加，分子为
$4A^2B^2+(A^2-B^2)^2=(A^2+B^2)^2$，所以总量为
$\omega^4(A^2+B^2)/2$，与直接求 $a_x^2+a_y^2$相同。只有先确认 $a_\tau=0$，才能用总加速度代替法向加速度；轴端点的简便算法不应直接套到整条椭圆上。
\end{drillsolution}

\end{document}
